\chapter{最优控制基础：变分、极小值原理、动态规划与 LQR}
\label{ch:optimal-control-basics}

% ════════════════════════════════════════════════════════════════
%  P4 控制部「C0 现代控制理论基础」子部 · 第 C0-5 章
%  主源：刘豹《现代控制理论》(第3版) 第6章（LP 引入·静态最优化·离散最优控制·
%        变分法→Euler-Lagrange→横截条件→极小值原理·DP/HJB·LQR物理意义+调节器/
%        跟踪器·Bang-Bang 双积分相平面）
%  辅源/第二信源：张嗣瀛/高立群《现代控制理论》(第2版) 第7章（四类端点变分解法·
%        横截条件·LQ Hamilton→Riccati 解析）
%  与导论 ch:control 的密度梯度去重：导论已高密度给 LQR 双推导/DARE/CARE/
%        Hamilton 矩阵法/iLQR/MPC（\cref{sec:ctrl-lqr,...}），本章只做「第一次
%        系统讲清」+ 刘豹独有的二次型物理意义逐项解析 + 跟踪器；严格化交 A/B/C 章。
%  OCR 勘误已逐式核对（刘豹 M2-d/e/f）：例6-5 e^{i}→e^t；例6-18 x(0)=[1;0]、
%        J*=2x1^2+2x1x2+x2^2；式6.310 \dot x=+dH/dlambda；式6.319 缺 \dot P、去 \times；
%        式6.316 末尾 \dot x→x；例6-20 β=√(q1/q2+a^2)（去导数点）；例6-22 +u→+u^2 等。
% ════════════════════════════════════════════════════════════════

\begin{practice}[前置自测]
开始之前，先试答下列问题。若已能流畅作答，可只速读本章的几何图解与速查卡；若卡住，正是本章要为你解决的。
\begin{enumerate}
  \item 求一个普通多元函数 $f(\mathbf{u})$ 的极小值，你会令 $\partial f/\partial\mathbf{u}=\mathbf{0}$。可如果待优化的不是一个\emph{向量}，而是一整条\emph{轨迹} $\mathbf{u}(t),\ t\in[t_0,t_f]$——“对一条函数求导”是什么意思？
  \item 把机械臂从一个位形在固定时间内挪到另一个位形，且“耗能最小”。这个“最优控制”问题里，谁是优化变量、谁是约束、目标泛函长什么样？
  \item 电机力矩有上限 $|u|\le u_{\max}$。当“最省时间”要求电机始终全力工作、只在某一瞬间反向时，为什么 $\partial H/\partial u=0$ 这一套\emph{失效}了？该用什么代替？
  \item 二次型代价 $\int(\mathbf{x}^\top\mathbf{Q}\mathbf{x}+\mathbf{u}^\top\mathbf{R}\mathbf{u})\,\mathrm dt$ 里，$\mathbf{Q}$ 的某个对角元调大、$\mathbf{R}$ 调大，闭环行为分别会怎样变？这两个权重各自“买”的是什么？
  \item 导论 \cref{sec:ctrl-lqr} 用动态规划与 HJB 两条路推出了 LQR 的 Riccati 方程。这两条路、再加上“变分/协态”第三条路，本质上是不是同一件事？它们从哪个更一般的框架派生？
\end{enumerate}
\end{practice}

\section*{本章目标}
学完本章，你应当能够：
\begin{enumerate}
  \item \textbf{区分}静态最优化（对向量求极值）与动态最优化（对轨迹/泛函求极值），并说清最优控制为何属于后者；
  \item \textbf{推导}泛函的变分、Euler--Lagrange 方程与各类\emph{横截条件}（固定/自由/可变端点），并用它求解无约束与状态方程约束下的最优轨迹；
  \item \textbf{建立}Hamilton 函数与协态（伴随）方程，把“带状态方程约束的泛函极值”系统地化为两点边值问题；
  \item \textbf{陈述并运用}Pontryagin 极小值原理，理解它在控制受界 $\mathbf{u}\in U$ 时如何\emph{替换}失效的平稳条件 $\partial H/\partial\mathbf{u}=\mathbf{0}$；
  \item \textbf{完整求解}最短时间（Bang--Bang）控制：从极小值原理导出 $\mathbf{u}^*=-\operatorname{SGN}(\mathbf{B}^\top\boldsymbol\lambda)$，并在双积分系统上画出相平面抛物线族与\emph{开关曲线}，给出状态反馈实现；
  \item \textbf{掌握}动态规划的最优性原理，从离散 Bellman 递推过渡到连续的 Hamilton--Jacobi--Bellman (HJB) 偏微分方程，并理解 DP 与极小值原理互为表里；
  \item \textbf{逐项解析}二次型代价各项的物理意义（误差代价/控制能量/终端代价），从而\emph{会选}权重 $\mathbf{Q},\mathbf{R},\mathbf{Q}_f$；
  \item \textbf{推导}有限/无限时域状态调节器的 Riccati 微分/代数方程与最优反馈 $\mathbf{u}^*=-\mathbf{K}\mathbf{x}$，并把它与导论的 LQR 结论焊接；
  \item \textbf{扩展}到输出调节器与\emph{跟踪器}：理解跟踪器为何多出一个前馈项 $\mathbf{g}(t)$、为何其现时控制依赖参考的\emph{全部未来值}；
  \item \textbf{联系}本章四种工具（变分/极小值原理/DP/LQR）的统一性，并知道每件工具在 A--Q 高级章里被严格化或推广到何处。
\end{enumerate}

\subsection*{本章知识导航}
本章是一条“\emph{从对向量求极值，逐步走到对轨迹求极值，再落到工程可实现的反馈律}”的主线。四件工具——\textbf{变分法、极小值原理、动态规划、LQR}——并非四套互不相干的方法，而是同一最优控制问题在不同假设下的四个侧面：变分法处理无约束情形并引出 Hamilton 框架；极小值原理把它推广到控制受界；动态规划从“分段决策”给出全局视角并落成 HJB；LQR 则是“线性系统 + 二次代价”这一最重要特例下三条路殊途同归的闭式解。结构如下：

\begin{center}
\small
\begin{tikzpicture}[
  node distance=5mm and 7mm, >=Stealth,
  blk/.style={draw, rounded corners, align=center, font=\small, inner sep=3pt, fill=main!6, draw=main!60!black},
  grp/.style={draw, rounded corners, align=center, font=\small, inner sep=3pt, fill=second!6, draw=second!60!black},
  ln/.style={->, thick, gray!60!black}]
\node[blk] (static) {静态最优化\\ $\partial f/\partial\mathbf{u}=\mathbf{0}$、Lagrange 乘子};
\node[blk, right=of static] (var) {变分法\\ Euler--Lagrange、横截条件};
\node[grp, right=of var] (ham) {Hamilton 框架\\ 协态 $\boldsymbol\lambda$、正则方程};
\node[grp, below=7mm of static] (pmp) {极小值原理\\ $\min_{\mathbf{u}\in U}H$};
\node[grp, right=of pmp] (bb) {Bang--Bang\\ 开关曲线、相平面};
\node[blk, right=of bb, fill=third!8, draw=third!60!black] (dp) {动态规划\\ Bellman $\to$ HJB};
\node[blk, below=7mm of pmp, fill=third!8, draw=third!60!black] (lqr) {LQR\\ 二次型物理意义};
\node[blk, right=of lqr, fill=third!8, draw=third!60!black] (ric) {Riccati\\ 调节器 / 跟踪器};
\draw[ln] (static)--(var); \draw[ln] (var)--(ham);
\draw[ln] (ham.south) -- ++(0,-3mm) -| (pmp.north);
\draw[ln] (pmp)--(bb); \draw[ln] (bb)--(dp);
\draw[ln] (pmp.south) -- (lqr.north);
\draw[ln] (lqr)--(ric);
\draw[ln] (dp.south) to[out=-90,in=0] (ric.east);
\end{tikzpicture}
\end{center}

\noindent\textbf{推荐路径}：第 \ref{sec:oc-overview}--\ref{sec:oc-static} 节是地基（什么是最优控制、静态最优化回顾）；第 \ref{sec:oc-variation}--\ref{sec:oc-hamilton} 节是变分主干（任何读者必经），建立 Euler--Lagrange 与 Hamilton 框架；第 \ref{sec:oc-pmp}--\ref{sec:oc-bangbang} 节讲极小值原理与最短时间控制（相平面图解是直觉精华）；第 \ref{sec:oc-dp} 节动态规划/HJB 给出全局视角；第 \ref{sec:oc-lqr} 节落到 LQR——\emph{这一节与导论 \cref{sec:ctrl-lqr} 高度互补}：导论已给双推导与 DARE/CARE，本章只补“二次型各项物理意义 + 调节器/跟踪器”这一直觉与实用面，不重推。

\subsection*{前置知识桥接}
本章是 P4「现代控制理论基础」子部的收官章，直接坐落在前四章地基之上：状态空间模型 $\dot{\mathbf{x}}=\mathbf{A}\mathbf{x}+\mathbf{B}\mathbf{u}$ 与状态转移（\cref{ch:statespace-basics}）、能控性（无限时域 LQR 存在性的前提，\cref{ch:controllability-obsv}）、Lyapunov 方程 $\mathbf{A}^\top\mathbf{P}+\mathbf{P}\mathbf{A}=-\mathbf{Q}$（LQR 闭环稳定性证明的工具，\cref{ch:lyapunov-basics}）、状态反馈/极点配置（综合的另一条线，\cref{ch:linear-synthesis}）。在数学侧，本章用到 P1 的二次型最小化与凸性（\cref{ch:nlopt}）、P0 的机械系统动力学记号 $\mathbf{M}(\mathbf{q})\ddot{\mathbf{q}}+\mathbf{C}\dot{\mathbf{q}}+\mathbf{g}=\boldsymbol\tau$（\cref{ch:rigid_body}）。微积分、变分初步、线性常微分方程求解假定已具备，不再重证。

\subsection*{与导论的分工}
\cref{ch:control} 作为控制导论，在 \cref{sec:ctrl-lqr} 已用\emph{导论密度}把 LQR 全线覆盖：离散动态规划完整推导（\cref{thm:ctrl-lqr-dp}）、连续 HJB 推导与 CARE（\cref{eq:ctrl-care}）、Hamilton 矩阵法解 Riccati（\cref{eq:ctrl-hamilton}）、Kalman“最优 $\ne$ 稳定”反例（\cref{ex:ctrl-kalman}）、乃至 iLQR/DDP（\cref{sec:ctrl-ilqr-ddp}）与 MPC。本章\textbf{不重复}这些已推导透的结论，而是补上导论\emph{没有}系统讲的两层：其一，LQR 之\emph{上游}——变分法、极小值原理、动态规划这三件更一般的工具，导论只在协态视角处嵌了一段（\cref{eq:ctrl-hamilton-sol} 邻近），本章给“第一次系统讲清”；其二，LQR 之\emph{侧面}——二次型代价各项的物理意义、有限时域 Riccati 微分方程的工程含义、以及导论未涉的\emph{输出调节器与跟踪器}。一言以蔽之：导论给“会用 LQR、会推 Riccati”，本章给“会用变分/极小值原理/DP 这套更一般的语言，并把 LQR 放回它应有的位置上、把权重的物理意义讲透”。凡两书都讲透之处（DARE/CARE 的存在唯一性、Hamilton 矩阵法、最优 $\ne$ 稳定），本章只 \cref{sec:ctrl-lqr} 式地指向导论，严格证明留给 \cref{ch:lqr-lqg-riccati}。

\begin{note}
\textbf{本章记号与约定.}\;
状态 $\mathbf{x}(t)\in\mathbb{R}^n$、控制 $\mathbf{u}(t)\in\mathbb{R}^r$ 一律粗体（教材源 $x,u$ 非粗体，本书统一粗体化）。性能泛函记 $J$，被积\emph{代价密度}记 $L(\mathbf{x},\mathbf{u},t)$（教材的“拉格朗日型”），终端代价记 $\Phi[\mathbf{x}(t_f)]$。协态/伴随矢量记 $\boldsymbol\lambda(t)\in\mathbb{R}^n$，Hamilton 函数 $H=L+\boldsymbol\lambda^\top\mathbf{f}$。
二次型代价权重沿用导论与全书纪律：状态权 $\mathbf{Q}\succeq\mathbf{0}$、控制权 $\mathbf{R}\succ\mathbf{0}$、终端权 $\mathbf{Q}_f\succeq\mathbf{0}$；cost-to-go/Riccati 矩阵记 $\mathbf{P}$（教材的 $\mathbf{Q}_1,\mathbf{Q}_2,\mathbf{Q}_0,\mathbf{P}$ 一律改写为 $\mathbf{Q},\mathbf{R},\mathbf{Q}_f,\mathbf{P}$）。
\textbf{$\mathbf{P}/\mathbf{Q}/\mathbf{R}$ 双关纪律}：本章 $\mathbf{P}$ 是 cost-to-go 权重、$\mathbf{Q},\mathbf{R}$ 是代价权重；勿与 \cref{ch:eskf} 估计侧的协方差 $\boldsymbol\Sigma$、Kalman 增益 $\mathbf{L}$ 混淆（两者经“估计--控制对偶”相关，但符号语义不同）。
最优量上标星号 $\mathbf{u}^*,\mathbf{x}^*,J^*$。本章默认\emph{确定性}最优控制（无随机噪声）。
\end{note}

% ════════════════════════════════════════════════════════════════
\section{什么是最优控制：从一个分配问题说起}
\label{sec:oc-overview}

\paragraph{动机：人人都在“求最优”。}
做任何事，人总想用最合理的方案收最好的效果——这背后就藏着最优化问题\cite{liubao2006modern}。最优\emph{控制}是其中针对动态系统的一支：它要为一个机组、一台设备或一个生产过程选定控制规律，使某种性能在给定意义下达到最优。在机器人学里，这正是“如何让机械臂以最省力的方式完成一次抓取”“如何让足式机器人以最短时间跨过一步”这类问题的数学母体。

\paragraph{反面：先看一个不含时间的“静态”最优化。}
为把“最优化”这个抽象词落地，先看一个与时间无关的例子\cite{liubao2006modern}。甲、乙两仓库分别存水泥 $1500$、$1800$ 包，A、B、C 三工地分别需 $900$、$600$、$1200$ 包；从甲库运往 A/B/C 每包运费 $1/2/4$ 元，从乙库运往 A/B/C 为 $4/5/9$ 元。设甲库运往三地的包数为 $x_1,x_2,x_3$、乙库为 $x_4,x_5,x_6$，则总运费
\begin{equation}
  f(\mathbf{x})=x_1+2x_2+4x_3+4x_4+5x_5+9x_6
  \label{eq:oc-lp-cost}
\end{equation}
是优化变量 $\mathbf{x}=(x_1,\dots,x_6)^\top$ 的\emph{目标函数}。任务是在约束
\begin{equation}
  x_1+x_2+x_3\le1500,\quad x_4+x_5+x_6\le1800,\quad x_1+x_4=900,\ x_2+x_5=600,\ x_3+x_6=1200
  \label{eq:oc-lp-constraints}
\end{equation}
下取 $f$ 的最小。目标与约束皆为 $\mathbf{x}$ 的一次函数，这是\emph{线性规划}；注意甲库运费处处更廉，故应尽先发甲库水泥，前两个不等式约束遂化为等式约束。这个例子的关键特征是：\textbf{优化变量是一个有限维向量，与时间无关}。

\paragraph{核心：动态系统让“变量”变成“函数”。}
最优控制与上例的根本区别在于，受控对象是一个动态系统——所有变量都是\emph{时间的函数}。于是目标不再是普通函数，而是“时间函数的函数”，称为\emph{泛函}（\cref{def:oc-functional}）。典型的目标泛函形如
\begin{equation}
  J=\int_{t_0}^{t_f}L[\mathbf{x}(t),\mathbf{u}(t),t]\,\mathrm dt,
  \label{eq:oc-integral-J}
\end{equation}
基本约束是受控对象的状态方程
\begin{equation}
  \dot{\mathbf{x}}(t)=\mathbf{f}[\mathbf{x}(t),\mathbf{u}(t),t].
  \label{eq:oc-state-eq}
\end{equation}
最优控制问题就是：在状态方程 \cref{eq:oc-state-eq} 的约束下，寻求最优控制函数 $\mathbf{u}^*(t)$，使目标泛函 \cref{eq:oc-integral-J} 取极值\cite{liubao2006modern,zhang2017modern}。

\begin{insight}
静态最优化与动态最优化并非泾渭分明。若把时域 $[t_0,t_f]$ 切成许多小段、每段内近似把变量当常量，动态最优化就近似成了一串\emph{分段静态最优化}——这正是“离散时间最优控制”的由来，也预告了\emph{动态规划}（\cref{sec:oc-dp}）“化整为零、逐段决策”的核心思想。分段越细，近似越准；令段长 $\to0$ 即回到连续问题。求泛函极值的工具——变分法——正是“求无穷多个变量的函数极值”，是微分法求极值的无穷维推广。
\end{insight}

\paragraph{最优控制问题的五要素。}
一个确定性最优控制问题由五部分完整确定\cite{liubao2006modern}：
\begin{enumerate}
  \item \textbf{动态描述}（状态方程）：连续 $\dot{\mathbf{x}}=\mathbf{f}[\mathbf{x},\mathbf{u},t]$，或离散 $\mathbf{x}(k+1)=\mathbf{f}[\mathbf{x}(k),\mathbf{u}(k),k]$。
  \item \textbf{控制作用域}（容许控制）：物理上 $\mathbf{u}(t)$ 不能任意大（电压/力矩有界），故约束在闭集 $\mathbf{u}(t)\in U=\{\mathbf{u}\mid\boldsymbol\varphi_j(\mathbf{x},\mathbf{u})\le0\}$ 内。落在 $U$ 内的 $\mathbf{u}$ 称\emph{容许控制}。
  \item \textbf{始端条件}：$t_0$ 给定；$\mathbf{x}(t_0)$ 固定（固定始端）、任意（自由始端）或受约束 $\mathbf{x}(t_0)\in\Omega_0$（可变始端）。
  \item \textbf{终端条件}：$t_f$ 与 $\mathbf{x}(t_f)$ 可固定、自由或受目标集约束 $\mathbf{x}(t_f)\in\Omega_f=\{\mathbf{x}(t_f)\mid\boldsymbol\varphi_j[\mathbf{x}(t_f)]=0\}$。
  \item \textbf{性能指标}（目标泛函）。
\end{enumerate}

\begin{definition}[三类性能指标]\label{def:oc-perf-types}
设连续系统，性能指标的一般形式（\emph{综合型}，又称 Bolza 型）为
\begin{equation}
  J=\Phi[\mathbf{x}(t_f)]+\int_{t_0}^{t_f}L[\mathbf{x}(t),\mathbf{u}(t),t]\,\mathrm dt,
  \label{eq:oc-bolza}
\end{equation}
其中 $\Phi$ 为\emph{终端指标函数}（反映对终态的要求，如脱靶量），$L$ 为\emph{动态指标函数}（反映过程品质、能量/燃料消耗）。两个退化情形：
\begin{itemize}
  \item $\Phi\equiv0$：$J=\int_{t_0}^{t_f}L\,\mathrm dt$，称\emph{积分型}（Lagrange 型）；
  \item $L\equiv0$：$J=\Phi[\mathbf{x}(t_f)]$，称\emph{终端型}（Mayer 型）。
\end{itemize}
\end{definition}

\noindent 满足约束、使 $J$ 取最小（最大）的控制称最优控制 $\mathbf{u}^*(t)$，其下状态方程的解称最优轨线 $\mathbf{x}^*(t)$，对应的 $J$ 值称最优指标 $J^*$。三大求解方法各擅胜场：\textbf{变分法}（\cref{sec:oc-variation}）处理无约束/光滑情形、\textbf{极小值原理}（\cref{sec:oc-pmp}）处理控制受界、\textbf{动态规划}（\cref{sec:oc-dp}）从全局/递推视角统摄二者。本书只讨论确定性问题；随机最优控制（与 \cref{ch:eskf} 的估计对偶）属另一支。

\begin{note}
\textbf{关于极小化约定.}\; 因 $f$ 取极小与 $-f$ 取极大完全等效，全章一律以\emph{极小化}陈述（教材亦然）。这与导论 LQR 的 $\min_U J$ 一致。
\end{note}

% ════════════════════════════════════════════════════════════════
\section{静态最优化：求泛函极值的“热身”}
\label{sec:oc-static}

\paragraph{动机：先把熟悉的函数极值理顺，触类旁通。}
动态最优化要求泛函极值，这要用变分法；而变分法与微分法在“求极值”上结构相似。故先回顾普通函数求极值，作为通向变分法的台阶\cite{liubao2006modern}。这一节的每个结论，到 \cref{sec:oc-variation} 都会有一个“函数 $\to$ 泛函”的对应版本。

\paragraph{无约束极值：梯度为零 + Hessian 正定。}
设 $n$ 元函数 $f=f(\mathbf{u})$，$\mathbf{u}=(u_1,\dots,u_n)^\top$。取极值的\emph{必要}条件是梯度为零矢量：
\begin{equation}
  \frac{\partial f}{\partial\mathbf{u}}=\mathbf{0},\quad\text{即}\quad\nabla f=\Big(\frac{\partial f}{\partial u_1},\dots,\frac{\partial f}{\partial u_n}\Big)^\top=\mathbf{0};
  \label{eq:oc-grad-zero}
\end{equation}
取\emph{极小}的充分条件还需 Hessian 矩阵正定：
\begin{equation}
  \frac{\partial^2 f}{\partial\mathbf{u}^2}=\Big[\frac{\partial^2 f}{\partial u_i\partial u_j}\Big]\succ\mathbf{0}.
  \label{eq:oc-hessian}
\end{equation}

\begin{example}[多元函数极值]\label{ex:oc-static-quad}
求 $f(\mathbf{x})=2x_1^2+5x_2^2+x_3^2+2x_2x_3+2x_3x_1-6x_2+3$ 的极小点。由 $\nabla f=\mathbf{0}$：$4x_1+2x_3=0$、$10x_2+2x_3-6=0$、$2x_3+2x_2+2x_1=0$，联立得 $\mathbf{x}^*=(1,1,-2)^\top$。Hessian
\[
  \frac{\partial^2 f}{\partial\mathbf{x}^2}=\begin{pmatrix}4&0&2\\0&10&2\\2&2&2\end{pmatrix}\succ\mathbf{0}
\]
正定，故 $\mathbf{x}^*$ 为极小点，$f^*=f(\mathbf{x}^*)=0$。
\end{example}

\paragraph{等式约束极值：Lagrange 乘子法。}
若极值受等式约束 $\mathbf{g}(\mathbf{x},\mathbf{u})=\mathbf{0}$，可用\emph{嵌入法}（从约束解出一个变量回代，仅适于简单情形）或更一般的\emph{Lagrange 乘子法}\cite{liubao2006modern}。后者引入乘子矢量 $\boldsymbol\lambda$，构造\emph{Lagrange 函数}（亦称 Hamilton 函数）
\begin{equation}
  H=J+\boldsymbol\lambda^\top\mathbf{g}=f(\mathbf{x},\mathbf{u})+\boldsymbol\lambda^\top\mathbf{g}(\mathbf{x},\mathbf{u}),
  \label{eq:oc-lagrangian}
\end{equation}
将带约束问题化为对 $H$ 的无约束极值问题。极值必要条件为
\begin{equation}
  \frac{\partial H}{\partial\mathbf{x}}=\mathbf{0},\qquad\frac{\partial H}{\partial\mathbf{u}}=\mathbf{0},\qquad\frac{\partial H}{\partial\boldsymbol\lambda}=\mathbf{g}(\mathbf{x},\mathbf{u})=\mathbf{0}.
  \label{eq:oc-lagrange-cond}
\end{equation}

\begin{example}[二次型 + 线性约束：调节器问题的“静态原型”]\label{ex:oc-static-lagrange}
求使 $J=\tfrac12\mathbf{x}^\top\mathbf{Q}\mathbf{x}+\tfrac12\mathbf{u}^\top\mathbf{R}\mathbf{u}$ 极小、满足约束 $\mathbf{g}(\mathbf{x},\mathbf{u})=\mathbf{x}+\mathbf{F}\mathbf{u}+\mathbf{d}=\mathbf{0}$ 的 $\mathbf{x}^*,\mathbf{u}^*$，其中 $\mathbf{Q},\mathbf{R}\succ\mathbf{0}$。构造 $H=\tfrac12\mathbf{x}^\top\mathbf{Q}\mathbf{x}+\tfrac12\mathbf{u}^\top\mathbf{R}\mathbf{u}+\boldsymbol\lambda^\top(\mathbf{x}+\mathbf{F}\mathbf{u}+\mathbf{d})$，由 \cref{eq:oc-lagrange-cond}：
\[
  \mathbf{Q}\mathbf{x}+\boldsymbol\lambda=\mathbf{0},\quad\mathbf{R}\mathbf{u}+\mathbf{F}^\top\boldsymbol\lambda=\mathbf{0},\quad\mathbf{x}+\mathbf{F}\mathbf{u}+\mathbf{d}=\mathbf{0}.
\]
联立解得
\begin{equation}
  \mathbf{u}^*=-(\mathbf{R}+\mathbf{F}^\top\mathbf{Q}\mathbf{F})^{-1}\mathbf{F}^\top\mathbf{Q}\,\mathbf{d}.
  \label{eq:oc-static-u}
\end{equation}
$\mathbf{Q},\mathbf{R}\succ\mathbf{0}$ 保证极小。\textbf{留意 \cref{eq:oc-static-u} 的结构}：$(\mathbf{R}+\mathbf{F}^\top\mathbf{Q}\mathbf{F})^{-1}\mathbf{F}^\top\mathbf{Q}$——它与 LQR 增益 $\mathbf{K}=(\mathbf{R}+\mathbf{B}^\top\mathbf{P}\mathbf{B})^{-1}\mathbf{B}^\top\mathbf{P}\mathbf{A}$（导论 \cref{eq:ctrl-dlqr-K}）形同骨肉。这绝非巧合：离散 LQR 不过是把这个静态二次型问题沿时间“叠 $N$ 层”，再用动态规划逐层求解（\cref{sec:oc-lqr-discrete-preview}）。
\end{example}

\begin{insight}
\cref{ex:oc-static-lagrange} 是理解整个最优控制的“种子”。\textbf{二次代价 + 线性约束 $\Rightarrow$ 线性最优解}——这条规律在静态、离散、连续三个层面反复重演，最终结晶为 LQR 的“线性反馈 + 二次 cost-to-go”。把这个 $6\times6$ 静态例子吃透，后面所有 Riccati 都只是它的“时间展开版”。
\end{insight}

% ════════════════════════════════════════════════════════════════
\section{变分法：对一整条轨迹求极值}
\label{sec:oc-variation}

\paragraph{动机：当“变量”有无穷多个。}
离散情形把 $[t_0,t_f]$ 分成有限段，变量数有限，可在有限维空间用静态方法处理。但令段长 $\to0$，变量成了 $t$ 的连续函数——无穷多个变量。从“求有限个变量的函数极值”到“求无穷个变量的泛函极值”，正是\emph{变分法}的用武之地\cite{liubao2006modern,zhang2017modern}。

\subsection{泛函、变分与泛函极值定理}
\label{sec:oc-functional}

\begin{definition}[泛函]\label{def:oc-functional}
若对某一类函数中每个确定的函数 $y(\cdot)$（\emph{整个函数}，而非某点函数值），因变量 $J$ 都有一确定的值与之对应，则称 $J$ 为宗量函数 $y(\cdot)$ 的\emph{泛函}，记 $J=J[y(\cdot)]$。简言之：泛函是“函数的函数”，其宗量是函数，而普通函数的宗量是变数。
\end{definition}

\begin{example}[最短弧长：变分问题的原型]\label{ex:oc-arclength}
平面上两点 $A(x_a,y_a)$、$B(x_b,y_b)$ 间，沿曲线 $y=y(x)$ 的弧长是该曲线的泛函。由弧长微分 $(\mathrm dl)^2=(\mathrm dx)^2+(\mathrm dy)^2$ 得 $\mathrm dl/\mathrm dx=\sqrt{1+\dot y^2}$（记 $\dot y=\mathrm dy/\mathrm dx$），故
\begin{equation}
  J[y(x)]=\int_{x_a}^{x_b}\sqrt{1+\dot y^2}\,\mathrm dx=\int_{x_a}^{x_b}L(\dot y)\,\mathrm dx.
  \label{eq:oc-arclength}
\end{equation}
不同曲线给不同弧长；显然两点间最短弧长是直线。这个问题在 \cref{ex:oc-intercept} 会作为“拦截/最短距离”再现。
\end{example}

\begin{figure}[ht]
\centering
\begin{tikzpicture}[scale=1.1,>=Stealth]
  \draw[->] (-0.3,0)--(4.3,0) node[right]{$x$};
  \draw[->] (0,-0.3)--(0,2.7) node[above]{$y$};
  \fill (0.5,0.6) circle (1.4pt) node[left]{$A$};
  \fill (3.7,2.2) circle (1.4pt) node[right]{$B$};
  % several candidate curves
  \draw[second, thick] (0.5,0.6)--(3.7,2.2) node[midway,below right,font=\scriptsize]{$y^*$ (最短)};
  \draw[main!70!black, thick] (0.5,0.6) .. controls (1.6,2.0) and (2.8,1.0) .. (3.7,2.2);
  \draw[main!70!black, thick] (0.5,0.6) .. controls (1.8,0.2) and (2.6,2.0) .. (3.7,2.2);
  \node[main!70!black,font=\scriptsize] at (2.2,2.05) {容许曲线 $y(x)$};
\end{tikzpicture}
\caption{弧长泛函 \cref{eq:oc-arclength}：每条连接 $A,B$ 的曲线 $y(x)$ 对应一个弧长值 $J[y(x)]$，最短者 $y^*$ 是直线。变分法要在所有“邻近”曲线中找出使 $J$ 极小的那条。}
\label{fig:oc-arclength}
\end{figure}

\paragraph{泛函的增量与变分。}
泛函的增量可表为线性主部加高阶项：
\begin{equation}
  \Delta J=J[y(x)+\delta y(x)]-J[y(x)]=\underbrace{L[y,\delta y]}_{\text{线性主部}}+\underbrace{R[y,\delta y]}_{\text{高阶项}},
  \label{eq:oc-increment}
\end{equation}
其中 $\delta y(x)$ 为宗量的\emph{变分}。定义增量的线性主部为泛函的\emph{变分} $\delta J:=L[y,\delta y]$。一个等价且便于计算的定义是
\begin{equation}
  \delta J=\frac{\partial}{\partial a}J[y(x)+a\,\delta y(x)]\Big|_{a=0}.
  \label{eq:oc-variation-def}
\end{equation}
\cref{eq:oc-variation-def} 把“对泛函求变分”化成“对一个实参数 $a$ 求导”，从而可用普通微分法则操作。一个反复用到的性质是“变分与求导可交换”：$\dfrac{\mathrm d}{\mathrm dx}\delta y=\delta\dot y$。

\begin{theorem}[泛函极值的必要条件]\label{thm:oc-variation-zero}
若可微泛函 $J[y(x)]$ 在 $y_0(x)$ 上达到极值，则其变分在 $y_0$ 上为零：$\delta J=0$。
\end{theorem}
\begin{proof}
设宗量变分 $\delta y$ 任意取定，则 $\varphi(a):=J[y_0(x)+a\,\delta y(x)]$ 是实变量 $a$ 的普通函数。由 \cref{eq:oc-variation-def}，$\dot\varphi(0)=\delta J[y_0]$。另一方面 $\varphi(0)=J[y_0]$ 已知是极值，由普通函数极值定理 $\dot\varphi(0)=0$。两者相等即得 $\delta J[y_0]=0$。
\end{proof}

\noindent 这正是 \cref{eq:oc-grad-zero}“梯度为零”的无穷维翻版——把“梯度”换成“变分”，把“等于零矢量”换成“变分泛函恒为零”。

\subsection{Euler--Lagrange 方程与横截条件}
\label{sec:oc-euler}

\paragraph{招牌推导：从 $\delta J=0$ 挤出 Euler--Lagrange 方程。}
考虑端点固定的积分泛函 $J=\int_{t_0}^{t_f}L(x,\dot x,t)\,\mathrm dt$，始点 $x(t_0)=x_0$、终点 $x(t_f)=x_f$。设极值曲线 $x^*(t)$，邻近曲线 $x(t)=x^*(t)+\varepsilon\eta(t)$，其中 $\eta(t)$ 任给连续可微、$\varepsilon$ 为小参数（\cref{fig:oc-euler}）。代入泛函，$J$ 成了 $\varepsilon$ 的函数，极值要求 $\partial J/\partial\varepsilon|_{\varepsilon=0}=0$：
\begin{equation}
  \int_{t_0}^{t_f}\Big[\eta(t)\frac{\partial L}{\partial x}+\dot\eta(t)\frac{\partial L}{\partial\dot x}\Big]\mathrm dt=0.
  \label{eq:oc-euler-step1}
\end{equation}
对第二项分部积分，$\int\dot\eta\,L_{\dot x}\,\mathrm dt=L_{\dot x}\eta\big|_{t_0}^{t_f}-\int\eta\,\tfrac{\mathrm d}{\mathrm dt}L_{\dot x}\,\mathrm dt$，代回得
\begin{equation}
  \int_{t_0}^{t_f}\eta(t)\Big[\frac{\partial L}{\partial x}-\frac{\mathrm d}{\mathrm dt}\frac{\partial L}{\partial\dot x}\Big]\mathrm dt+\frac{\partial L}{\partial\dot x}\eta(t)\Big|_{t_0}^{t_f}=0.
  \label{eq:oc-euler-step2}
\end{equation}
端点固定 $\Rightarrow\eta(t_0)=\eta(t_f)=0$，边界项消失。又因 $\eta(t)$ 在 $(t_0,t_f)$ 内任意，由变分法基本引理，方括号必恒为零：

\begin{theorem}[Euler--Lagrange 方程]\label{thm:oc-euler}
使 $J=\int_{t_0}^{t_f}L(x,\dot x,t)\,\mathrm dt$ 取极值的曲线 $x^*(t)$ 必满足二阶微分方程
\begin{equation}
  \boxed{\frac{\partial L}{\partial x}-\frac{\mathrm d}{\mathrm dt}\frac{\partial L}{\partial\dot x}=0,}
  \label{eq:oc-euler}
\end{equation}
其展开式为 $L_x-L_{\dot x t}-L_{\dot x x}\dot x-L_{\dot x\dot x}\ddot x=0$。其中 $x$ 须二阶连续可微、$L$ 至少两次连续可微。
\end{theorem}

\begin{figure}[ht]
\centering
\begin{tikzpicture}[scale=1.0,>=Stealth]
  \draw[->] (-0.3,0)--(5.3,0) node[right]{$t$};
  \draw[->] (0,-0.3)--(0,3.0) node[above]{$x$};
  \fill (0.5,0.7) circle (1.3pt) node[below left]{$x_0$};
  \fill (4.7,2.3) circle (1.3pt) node[below right]{$x_f$};
  \draw[xshift=0,dashed] (0.5,0)--(0.5,0.7); \node[font=\scriptsize] at (0.5,-0.25){$t_0$};
  \draw[dashed] (4.7,0)--(4.7,2.3); \node[font=\scriptsize] at (4.7,-0.25){$t_f$};
  \draw[main, very thick] (0.5,0.7) .. controls (2.0,2.0) and (3.2,1.6) .. (4.7,2.3)
     node[midway,above,font=\scriptsize,main]{$x^*(t)$ 极值曲线};
  \draw[second, thick] (0.5,0.7) .. controls (2.0,2.6) and (3.2,2.2) .. (4.7,2.3);
  \node[second,font=\scriptsize] at (2.9,2.55){$x^*(t)+\varepsilon\eta(t)$};
  % indicate eta
  \draw[<->,gray] (2.0,1.62) -- (2.0,2.28); \node[gray,font=\scriptsize] at (2.35,1.95){$\varepsilon\eta$};
\end{tikzpicture}
\caption{Euler--Lagrange 推导（端点固定）：在极值曲线 $x^*(t)$ 旁取一族邻近曲线 $x^*+\varepsilon\eta$，其中 $\eta(t_0)=\eta(t_f)=0$。令 $J$ 对 $\varepsilon$ 的导数在 $\varepsilon=0$ 处为零，分部积分后由 $\eta$ 的任意性逼出 \cref{eq:oc-euler}。（仿刘豹图 6.4 重绘）}
\label{fig:oc-euler}
\end{figure}

\begin{insight}
\textbf{Euler--Lagrange 是“梯度为零”的轨迹版。}\; 静态情形 $\partial f/\partial\mathbf{u}=\mathbf{0}$ 是一组\emph{代数}方程；泛函情形 $L_x-\tfrac{\mathrm d}{\mathrm dt}L_{\dot x}=0$ 是一个\emph{微分}方程。代数解给一个点，微分方程解给一条曲线。这就是“无穷维优化”的代价与馈赠：变量从一个点变成一条函数，最优性条件也从代数方程升级为微分方程。
\end{insight}

\paragraph{自由端点与横截条件。}
若回看 \cref{eq:oc-euler-step2}：当端点\emph{不}固定时，边界项 $\frac{\partial L}{\partial\dot x}\delta x\big|_{t_0}^{t_f}$ 不再自动为零，它必须独立地等于零——这给出补足边界条件的\emph{横截条件}：
\begin{equation}
  \frac{\partial L}{\partial\dot x}\Big|_{t_0}=0\ (\text{始端自由}),\qquad\frac{\partial L}{\partial\dot x}\Big|_{t_f}=0\ (\text{终端自由}).
  \label{eq:oc-transversality-free}
\end{equation}
Euler--Lagrange 是二阶微分方程，需两个边界条件：固定端点由 $x(t_0)=x_0,\ x(t_f)=x_f$ 提供；自由端点则由对应的横截条件补足。这构成一个\emph{两点边值问题}。

\begin{example}[二阶系统的最优轨迹]\label{ex:oc-euler-simple}
受控对象 $\dot x=u$，边界 $x(0)=x_0,\ x(t_f)=x_f$，求使 $J=\int_0^{t_f}(x^2+u^2)\,\mathrm dt$ 极小的 $u^*$。把 $u=\dot x$ 代入消去控制，$L=x^2+\dot x^2$。Euler--Lagrange 给 $2x-2\ddot x=0$，即 $\ddot x=x$，通解
\begin{equation}
  x(t)=C_1\mathrm e^{t}+C_2\mathrm e^{-t}.
  \label{eq:oc-euler-sol}
\end{equation}
\footnote{此处 $\mathrm e^{t}$ 而非教材 OCR 误印的 $\mathrm e^{i}$：通解为双曲型 $C_1\mathrm e^{t}+C_2\mathrm e^{-t}$，与极值曲线 $x^*=\frac{x_f\sinh t+x_0\sinh(t_f-t)}{\sinh t_f}$ 自洽（刘豹勘误 M2-e）。}
代入边界条件定 $C_1,C_2$，得极值曲线 $x^*(t)=\dfrac{x_f\sinh t+x_0\sinh(t_f-t)}{\sinh t_f}$，最优控制 $u^*=\dot x^*=\dfrac{x_f\cosh t-x_0\cosh(t_f-t)}{\sinh t_f}$。
\end{example}

\paragraph{可变端点与靶线横截条件。}
更一般地，若始端固定而终端时刻 $t_f$ 可沿给定\emph{靶线} $\mathbf x(t_f)=\mathbf C(t_f)$ 滑动（如火箭拦截沿轨道 $\mathbf C(t)$ 运动的目标，\cref{fig:oc-intercept}），则除 Euler--Lagrange 外，还要确定最优终端时刻 $t_f^*$。此时横截条件升级为\emph{靶线横截条件}\cite{liubao2006modern}：
\begin{equation}
  \Big\{L+[\dot{\mathbf C}(t)-\dot{\mathbf x}(t)]^\top\frac{\partial L}{\partial\dot{\mathbf x}}\Big\}\Big|_{t=t_f^*}=0.
  \label{eq:oc-transversality-target}
\end{equation}
两个常见退化：靶线平行 $t$ 轴（$\dot{\mathbf C}=\mathbf0$）时化为 $\{L-\dot{\mathbf x}^\top L_{\dot{\mathbf x}}\}|_{t_f}=0$；靶线垂直 $t$ 轴（$\dot{\mathbf C}=\infty$）时退回 $L_{\dot{\mathbf x}}|_{t_f}=0$，即 \cref{eq:oc-transversality-free}。

\begin{figure}[ht]
\centering
\begin{tikzpicture}[scale=1.0,>=Stealth]
  \draw[->] (-0.3,0)--(5.3,0) node[right]{$t$};
  \draw[->] (0,-0.3)--(0,3.0) node[above]{$x$};
  \fill (0.4,0.5) circle (1.3pt) node[below left]{$x(0)$};
  % target line C(t)
  \draw[second, very thick] (1.5,2.6) -- (4.6,0.6) node[right,font=\scriptsize,second]{靶线 $x=C(t)$};
  % optimal trajectory hitting target
  \draw[main, very thick] (0.4,0.5) .. controls (1.8,1.4) and (2.6,1.7) .. (3.35,1.65);
  \fill[main] (3.35,1.65) circle (1.4pt);
  \node[main,font=\scriptsize] at (1.7,1.05){$x^*(t)$};
  \draw[dashed,gray] (3.35,1.65)--(3.35,0); \node[font=\scriptsize] at (3.35,-0.25){$t_f^*$};
  \node[font=\scriptsize,align=center] at (4.0,2.0){交点处\\横截条件 \cref{eq:oc-transversality-target}};
\end{tikzpicture}
\caption{可变终端（拦截）问题：极值轨迹 $x^*(t)$ 须落在沿靶线 $C(t)$ 运动的目标上，终端时刻 $t_f^*$ 由靶线横截条件 \cref{eq:oc-transversality-target} 与终端约束 $x(t_f)=C(t_f)$ 共同确定。（仿刘豹图 6.6/6.7 重绘）}
\label{fig:oc-intercept}
\end{figure}

\begin{example}[最短距离到一条直线靶]\label{ex:oc-intercept}
求从 $x(0)=1$ 到直线 $x(t)=2-t$ 间距离最短的曲线。即极小化 $J=\int_0^{t_f}\sqrt{1+\dot x^2}\,\mathrm dt$，$L=\sqrt{1+\dot x^2}$。Euler--Lagrange 给 $\frac{\mathrm d}{\mathrm dt}\frac{\dot x}{\sqrt{1+\dot x^2}}=0\Rightarrow\dot x=\text{常数}=a$，故 $x=at+b$（直线，符合直觉）。由 $x(0)=1$ 得 $b=1$。靶线 $C(t)=2-t$，$\dot C=-1$，代入 \cref{eq:oc-transversality-target}：$-(1+\dot x)\frac{\dot x}{\sqrt{1+\dot x^2}}+\sqrt{1+\dot x^2}=0\Rightarrow\dot x=1$，故 $a=1$，最优轨线 $x^*=t+1$。代入终端约束 $t_f+1=2-t_f$ 解出 $t_f^*=\tfrac12$。
\textbf{几何意义}：$\dot x^*=1$ 与靶线斜率 $-1$ 之积为 $-1$，即最短路径\emph{垂直}于靶线——横截条件的几何本质就是“垂直相交”。
\end{example}

\paragraph{含终端项的综合型泛函。}
若泛函含终端项 $J=\Phi[\mathbf x(t_f)]+\int_{t_0}^{t_f}L\,\mathrm dt$，$t_f$ 给定、$\mathbf x(t_f)$ 自由，则仿上推导，终端横截条件被终端项“拉偏”：
\begin{equation}
  \frac{\partial L}{\partial\dot{\mathbf x}}\Big|_{t_f}=-\frac{\partial\Phi[\mathbf x(t_f)]}{\partial\mathbf x(t_f)}.
  \label{eq:oc-transversality-phi}
\end{equation}
这预告了 LQR 有限时域的终端条件 $\mathbf P(t_f)=\mathbf Q_f$（\cref{eq:oc-rde}）：终端权 $\mathbf Q_f$ 正是经此横截条件进入 Riccati 边界的。

\begin{pitfall}[只解 Euler--Lagrange、忘了横截条件]
Euler--Lagrange 是二阶微分方程，必须配\emph{两个}边界条件才能定下两个积分常数。固定端点提供两个；只要有一端自由，就缺一个条件，\textbf{必须}用对应的横截条件补足。新手常只写 Euler--Lagrange、用一个边界条件硬凑，结果或欠定或错解。记号上四类端点（固定/自由始端、固定/自由终端、可变 $t_f$）各有专属横截条件——这正是 \cref{ch:control} 在协态视角处一笔带过、而本章作为“做透层”逐条铺开的内容（第二信源见现控02 第7章四类端点解法\cite{zhang2017modern}）。
\end{pitfall}

% ════════════════════════════════════════════════════════════════
\section{Hamilton 框架：把状态方程约束并入泛函}
\label{sec:oc-hamilton}

\paragraph{动机：真实问题里 $x$ 与 $u$ 被状态方程绑定。}
\cref{ex:oc-euler-simple} 取巧地把 $u=\dot x$ 消去了控制。但一般状态方程 $\dot{\mathbf x}=\mathbf f[\mathbf x,\mathbf u,t]$ 无法这样消元——$\mathbf x$ 与 $\mathbf u$ 是被状态方程\emph{约束}的两组独立函数。解法与静态情形 \cref{eq:oc-lagrangian} 完全平行：用 Lagrange 乘子（这里是\emph{协态函数} $\boldsymbol\lambda(t)$）把约束并入泛函\cite{liubao2006modern,zhang2017modern}。

\paragraph{构造 Hamilton 函数与增广泛函。}
把状态方程写成约束 $\mathbf f[\mathbf x,\mathbf u,t]-\dot{\mathbf x}=\mathbf0$，引入 $n$ 维协态 $\boldsymbol\lambda(t)$，构造增广泛函
\begin{equation}
  J'=\int_{t_0}^{t_f}\big\{L[\mathbf x,\mathbf u,t]+\boldsymbol\lambda^\top(t)[\mathbf f[\mathbf x,\mathbf u,t]-\dot{\mathbf x}]\big\}\mathrm dt.
  \label{eq:oc-augmented}
\end{equation}
定义\emph{Hamilton 函数}
\begin{equation}
  \boxed{H[\mathbf x,\mathbf u,\boldsymbol\lambda,t]=L[\mathbf x,\mathbf u,t]+\boldsymbol\lambda^\top(t)\,\mathbf f[\mathbf x,\mathbf u,t],}
  \label{eq:oc-hamiltonian}
\end{equation}
则 $J'=\int_{t_0}^{t_f}\{H-\boldsymbol\lambda^\top\dot{\mathbf x}\}\mathrm dt$。对 $-\boldsymbol\lambda^\top\dot{\mathbf x}$ 项分部积分 $\int(-\boldsymbol\lambda^\top\dot{\mathbf x})\mathrm dt=\int\dot{\boldsymbol\lambda}^\top\mathbf x\,\mathrm dt-\boldsymbol\lambda^\top\mathbf x|_{t_0}^{t_f}$，得 $J'=\int_{t_0}^{t_f}\{H+\dot{\boldsymbol\lambda}^\top\mathbf x\}\mathrm dt-\boldsymbol\lambda^\top\mathbf x|_{t_0}^{t_f}$。计算 $\delta\mathbf u,\delta\mathbf x$ 引起的一阶变分并令 $\delta J'=0$（对任意 $\delta\mathbf u,\delta\mathbf x$），得必要条件。

\begin{theorem}[Lagrange 问题的必要条件：正则方程]\label{thm:oc-canonical}
对状态方程 $\dot{\mathbf x}=\mathbf f[\mathbf x,\mathbf u,t]$ 约束下的积分型泛函 $J=\int_{t_0}^{t_f}L\,\mathrm dt$（控制\emph{无界}），最优解满足
\begin{align}
  \dot{\mathbf x}&=\frac{\partial H}{\partial\boldsymbol\lambda}=\mathbf f[\mathbf x,\mathbf u,t] &&\text{(状态方程)},\label{eq:oc-canonical-x}\\
  \dot{\boldsymbol\lambda}&=-\frac{\partial H}{\partial\mathbf x} &&\text{(协态/伴随方程)},\label{eq:oc-canonical-lambda}\\
  \mathbf0&=\frac{\partial H}{\partial\mathbf u} &&\text{(控制方程/平稳条件)},\label{eq:oc-canonical-u}
\end{align}
配横截条件 $\boldsymbol\lambda|_{t_0}^{t_f}$ 给出边界。\cref{eq:oc-canonical-x,eq:oc-canonical-lambda} 合称\emph{Hamilton 正则方程}。$\boldsymbol\lambda$ 称协态/伴随矢量。
\end{theorem}

\noindent 边界条件随端点而定：始端固定 + 终态自由 $\Rightarrow\mathbf x(t_0)=\mathbf x_0,\ \boldsymbol\lambda(t_f)=\mathbf0$；两端固定 $\Rightarrow\mathbf x(t_0)=\mathbf x_0,\ \mathbf x(t_f)=\mathbf x_f$。求解三步：(1) 由控制方程 $\partial H/\partial\mathbf u=\mathbf0$ 解出 $\mathbf u^*=\tilde{\mathbf u}[\mathbf x,\boldsymbol\lambda]$；(2) 代入正则方程解两点边值问题求 $\mathbf x^*,\boldsymbol\lambda^*$；(3) 回代得 $\mathbf u^*$。

\begin{insight}
\textbf{协态 $\boldsymbol\lambda$ 是“影子价格”。}\; 静态 Lagrange 乘子衡量“约束松一点能省多少代价”；动态的协态 $\boldsymbol\lambda(t)$ 同理——它是“在 $t$ 时刻把状态推动一点点，对剩余最优代价的边际影响”。事实上 \cref{sec:oc-dp} 会证明 $\boldsymbol\lambda(t)=\partial J^*/\partial\mathbf x$（值函数对状态的梯度）。这把变分（协态）、动态规划（值函数梯度）两套语言焊在一起。导论 \cref{eq:ctrl-hamilton} 用的 Hamilton 矩阵正是把这套正则方程线性化、写成 $2n\times2n$ 矩阵的结果。
\end{insight}

\paragraph{Hamilton 函数沿最优轨线的守恒。}
$H$ 对时间的全导数为 $\frac{\mathrm dH}{\mathrm dt}=\frac{\partial H}{\partial t}+(\frac{\partial H}{\partial\mathbf u})^\top\dot{\mathbf u}+(\frac{\partial H}{\partial\mathbf x}+\dot{\boldsymbol\lambda})^\top\mathbf f$。沿最优轨线 $\partial H/\partial\mathbf u=\mathbf0$、$\partial H/\partial\mathbf x+\dot{\boldsymbol\lambda}=\mathbf0$，故
\begin{equation}
  \frac{\mathrm dH}{\mathrm dt}=\frac{\partial H}{\partial t}.
  \label{eq:oc-H-conserved}
\end{equation}
当 $H$ 不显含 $t$（定常系统 + 定常代价）时 $\mathrm dH/\mathrm dt=0$，即\textbf{沿最优轨线 $H$ 恒为常值}——这是力学中“Hamilton 量守恒”在最优控制里的影子。

\begin{example}[两端固定的最省能量控制]\label{ex:oc-double-int-fixed}
双积分系统 $\dot x_1=x_2,\ \dot x_2=u$（如惯性负载的位置-速度），要在 $2\,\mathrm s$ 内从 $\mathbf x(0)=(1,1)^\top$ 转移到 $\mathbf x(2)=(0,0)^\top$，极小化 $J=\tfrac12\int_0^2u^2\,\mathrm dt$（省能量）。Hamilton 函数 $H=\tfrac12u^2+\lambda_1x_2+\lambda_2u$。控制方程 $\partial H/\partial u=u+\lambda_2=0\Rightarrow u=-\lambda_2$；协态方程 $\dot\lambda_1=0,\ \dot\lambda_2=-\lambda_1$，解得 $\lambda_1=C_1,\ \lambda_2=-C_1t+C_2$，故 $u=C_1t-C_2$。积分状态方程并由四个边界条件 $x_1(0)=1,x_2(0)=1,x_1(2)=0,x_2(2)=0$ 定 $C_1=3,C_2=\tfrac72,C_3=1,C_4=1$，得
\begin{equation}
  u^*(t)=3t-\tfrac72,\quad x_1^*(t)=\tfrac12t^3-\tfrac74t^2+t+1,\quad x_2^*(t)=\tfrac32t^2-\tfrac74t+1.
  \label{eq:oc-double-int-sol}
\end{equation}
注意 $u^*(t)$ 是 $t$ 的\emph{线性}函数（二次代价 + 线性系统 $\Rightarrow$ 协态线性 $\Rightarrow$ 控制线性）——又一次印证 \cref{ex:oc-static-lagrange} 的“种子”规律。若改终端为部分自由（$x_2(2)$ 自由），则用横截条件 $\lambda_2(2)=0$ 替换 $x_2(2)=0$，解变为 $u^*(t)=\tfrac98t-\tfrac94\cdot\tfrac{?}{}$\footnote{改 $x_2(2)$ 自由后 $C_1=\tfrac98,C_2=\tfrac{18}{8}$，$u^*(t)=\tfrac98t-\tfrac94$（刘豹勘误 M2-d：教材 OCR 误印 $6t-12$，由 $x_2^*=\tfrac{9}{16}t^2-\tfrac94t+1$ 求导即得 $\dot x_2^*=\tfrac98t-\tfrac94=u^*$）。}——同一系统、不同端点条件给出不同最优解，凸显了横截条件的分量。
\end{example}

\paragraph{Bolza 问题：终端项 + 终端约束。}
最一般的（Bolza）情形：$J=\Phi[\mathbf x(t_f),t_f]+\int_{t_0}^{t_f}L\,\mathrm dt$，终态满足约束 $\mathbf N[\mathbf x(t_f),t_f]=\mathbf0$，$t_f$ 待定。需引入两个乘子：$n$ 维 $\boldsymbol\lambda(t)$（约束状态方程）与 $q$ 维 $\boldsymbol\mu$（约束终端流形）。可变端点变分给出必要条件：正则方程 \cref{eq:oc-canonical-x,eq:oc-canonical-lambda,eq:oc-canonical-u} 不变，但终端条件变为
\begin{equation}
  \boldsymbol\lambda(t_f)=\frac{\partial\Phi}{\partial\mathbf x(t_f)}+\frac{\partial\mathbf N^\top}{\partial\mathbf x(t_f)}\boldsymbol\mu,\qquad
  \Big[H+\frac{\partial\Phi}{\partial t_f}+\frac{\partial\mathbf N^\top}{\partial t_f}\boldsymbol\mu\Big]\Big|_{t_f}=0.
  \label{eq:oc-bolza-terminal}
\end{equation}
第二式用于确定自由的 $t_f$（终端时刻固定时该式消失）。这套“$2n+r+q+1$ 个方程定 $2n+r+q+1$ 个未知量”的框架，是变分法处理一般最优控制问题的终点；其图解（拦截、可变终端各变分关系）见 \cref{fig:oc-intercept}。

\begin{example}[终态受约束的最省能量]\label{ex:oc-bolza}
双积分系统 $\dot x_1=x_2,\dot x_2=u$，$\mathbf x(0)=\mathbf0$，终态约束 $x_1(1)+x_2(1)-1=0$（一条直线流形），$t_f=1$，极小化 $J=\tfrac12\int_0^1u^2\,\mathrm dt$。$H=\tfrac12u^2+\lambda_1x_2+\lambda_2u$，正则方程同 \cref{ex:oc-double-int-fixed}。终端条件由 \cref{eq:oc-bolza-terminal} 给 $\lambda_1(1)=\mu\partial N/\partial x_1=\mu$、$\lambda_2(1)=\mu$。代入通解定常数，再由终端约束 $x_1(1)+x_2(1)=1$ 解出 $\mu=-\tfrac37$，得 $u^*(t)=-\tfrac37t+\tfrac67$，$x_1^*=-\tfrac1{14}t^3+\tfrac37t^2$，$x_2^*=-\tfrac3{14}t^2+\tfrac67t$。
\end{example}

% ════════════════════════════════════════════════════════════════
\section{Pontryagin 极小值原理：当控制受界}
\label{sec:oc-pmp}

\paragraph{动机：电机力矩有上限，平稳条件失效。}
变分法的控制方程 $\partial H/\partial\mathbf u=\mathbf0$ 默认 $\mathbf u$ 可任意取值、$\delta\mathbf u$ 任意。但工程中控制几乎总是受界的：力矩、电压、推力都有上限，$\mathbf u(t)\in U$（有界闭集），$\delta\mathbf u$ 不能任意。这时 $\partial H/\partial\mathbf u=\mathbf0$ \emph{不成立}——正如定义在闭区间上的函数不能靠“导数为零”判断端点处的极值\cite{liubao2006modern}。Pontryagin（庞特里亚金）1956 年提出的\emph{极小值原理}恰好补上这一缺口\cite{pontryagin1962mathematical}。

\begin{theorem}[Pontryagin 极小值原理]\label{thm:oc-pmp}
设系统 $\dot{\mathbf x}=\mathbf f[\mathbf x,\mathbf u,t],\ \mathbf x(t_0)=\mathbf x_0$，控制受约束 $\mathbf u\in U$，终端约束 $\mathbf N[\mathbf x(t_f),t_f]=\mathbf0$（$t_f$ 可待定），性能泛函 $J=\Phi[\mathbf x(t_f),t_f]+\int_{t_0}^{t_f}L\,\mathrm dt$，Hamilton 函数 $H=L+\boldsymbol\lambda^\top\mathbf f$。则实现最优控制的\emph{必要}条件是：
\begin{enumerate}
  \item 沿最优轨线满足正则方程 $\dot{\mathbf x}=\partial H/\partial\boldsymbol\lambda$，$\dot{\boldsymbol\lambda}=-\partial H/\partial\mathbf x$（$\mathbf g$ 不含 $\mathbf x$ 时）；
  \item \textbf{在最优轨线上，最优控制 $\mathbf u^*$ 使 $H$ 取\emph{全局极小}}：
  \begin{equation}
    \boxed{H[\mathbf x^*,\boldsymbol\lambda^*,\mathbf u^*,t]=\min_{\mathbf u\in U}H[\mathbf x^*,\boldsymbol\lambda^*,\mathbf u,t];}
    \label{eq:oc-pmp-min}
  \end{equation}
  \item $H$ 在终点的值由 $[H+\partial\Phi/\partial t_f+\boldsymbol\mu^\top\partial\mathbf N/\partial t_f]|_{t_f}=0$ 确定（$t_f$ 自由时）；
  \item 协态终值满足横截条件 \cref{eq:oc-bolza-terminal}，配边界 $\mathbf x(t_0)=\mathbf x_0,\ \mathbf N[\mathbf x(t_f),t_f]=\mathbf0$。
\end{enumerate}
\end{theorem}

\paragraph{关键替换：从“平稳”到“全局极小”。}
极小值原理与变分法必要条件的唯一区别，就是把平稳条件 $\partial H/\partial\mathbf u=\mathbf0$ \emph{替换}为全局极小 \cref{eq:oc-pmp-min}（横截条件、端点条件不变）。这一替换的威力在于：当 $H$ 对 $\mathbf u$ 是\emph{线性}函数，或在 $U$ 内\emph{单调}时，$\partial H/\partial\mathbf u\ne\mathbf0$，变分法束手无策，而极小值原理直接告诉你最优控制取\emph{边界值}。可证：对线性系统，极小值原理既必要也充分。

\begin{insight}
\textbf{“极小值”一词的来历与“极大值原理”的别名。}\; 名字来自结论 \cref{eq:oc-pmp-min}——最优控制使 $H$ 全局\emph{极小}。若把 $\boldsymbol\lambda$ 与 $H$ 的符号反过来（$\bar H=-H$），同一定理就成了“$\bar H$ 取全局极大”，故有些文献（尤其俄文原典与早期英译）称\emph{极大值原理}（maximum principle）\cite{pontryagin1962mathematical}。二者只差一个符号约定，内容完全等同。本书与刘豹一致用“极小值原理”。
\end{insight}

\begin{example}[控制线性 $\Rightarrow$ 取边界值]\label{ex:oc-pmp-linear}
系统 $\dot x=x-u,\ x(0)=5$，控制约束 $\tfrac12\le u\le1$，极小化 $J=\int_0^1(x+u)\,\mathrm dt$。$H=x+u+\lambda(x-u)=x(1+\lambda)+u(1-\lambda)$ 对 $u$ \emph{线性}，$\partial H/\partial u=1-\lambda$ 与 $u$ 无关——平稳条件失效。由极小值原理，欲使 $H$ 极小只需使 $u(1-\lambda)$ 极小：$\lambda>1$ 取上界 $u^*=1$，$\lambda<1$ 取下界 $u^*=\tfrac12$。协态方程 $\dot\lambda=-(1+\lambda)$，配 $\lambda(1)=0$（终态自由）解得 $\lambda=\mathrm e^{1-t}-1$，令 $\lambda=1$ 得切换时刻 $t=1-\ln2\approx0.307$。于是 $u^*$ 在 $t<0.307$ 取 $1$、$t>0.307$ 取 $\tfrac12$——\textbf{控制在边界间切换，这是变分法绝对求不出来的解}。
\end{example}

% ════════════════════════════════════════════════════════════════
\section{最短时间与 Bang--Bang 控制：相平面图解}
\label{sec:oc-bangbang}

\paragraph{动机：最快是怎样炼成的。}
极小值原理最漂亮的应用是\emph{最短时间控制}。它的性能指标极简（$L=1$ 或 $\Phi=t_f,L=0$），因而研究最早、结论最熟，也最能展示极小值原理的几何威力。设能控线性定常系统 $\dot{\mathbf x}=\mathbf A\mathbf x+\mathbf B\mathbf u$，$\mathbf x(0)=\mathbf x_0,\ \mathbf x(t_f)=\mathbf0$（$t_f$ 待求），控制约束 $|u_i|\le1$，极小化 $J=\int_0^{t_f}1\,\mathrm dt=t_f$\cite{liubao2006modern}。

\begin{theorem}[最短时间控制是 Bang--Bang]\label{thm:oc-bangbang}
上述最短时间问题的 Hamilton 函数 $H=1+\boldsymbol\lambda^\top(\mathbf A\mathbf x+\mathbf B\mathbf u)=1+\mathbf x^\top\mathbf A^\top\boldsymbol\lambda+\mathbf u^\top\mathbf B^\top\boldsymbol\lambda$。为使 $H$ 全局极小（$\mathbf u^\top\mathbf B^\top\boldsymbol\lambda$ 极小、$|u_i|\le1$），最优控制取\emph{符号反向的边界值}：
\begin{equation}
  \boxed{\mathbf u^*(t)=-\operatorname{SGN}[\mathbf B^\top\boldsymbol\lambda(t)],\qquad u_i^*(t)=-\operatorname{sgn}[(\mathbf B^\top\boldsymbol\lambda)_i].}
  \label{eq:oc-bangbang-u}
\end{equation}
由协态方程 $\dot{\boldsymbol\lambda}=-\mathbf A^\top\boldsymbol\lambda$ 解得 $\boldsymbol\lambda(t)=\mathrm e^{-\mathbf A^\top t}\boldsymbol\lambda_0$，故 $\mathbf u^*(t)=-\operatorname{SGN}[\mathbf B^\top\mathrm e^{-\mathbf A^\top t}\boldsymbol\lambda_0]$。控制只取边界值 $\pm1$、并随 $\boldsymbol\lambda$ 变号而切换——这种“开关控制”称\emph{Bang--Bang 控制}。
\end{theorem}

\noindent 注意 $\boldsymbol\lambda_0\ne\mathbf0$：否则 $\boldsymbol\lambda\equiv\mathbf0$，由 $H$ 表达式将导出 $1=0$ 的矛盾。两个经典结论：(i) 若存在最短时间控制，则它\emph{唯一}（反证：若有两个不同的最优控制 $\mathbf u^*\ne\mathbf v^*$，其平均 $\mathbf w=\tfrac12(\mathbf u^*+\mathbf v^*)$ 也最优，但在两者不等处 $\mathbf w$ 取内点值、非 Bang--Bang，矛盾）；(ii) 若 $\mathbf A$ 特征值\emph{全为实数}，则每个 $u_i$ 至多切换 $n-1$ 次（$\mathbf A$ 有复特征值时切换次数无上界）。

\subsection{双积分系统的开关曲线}
\label{sec:oc-double-int-bb}

\paragraph{从极小值原理到相平面。}
取最简的双积分系统 $\dot x_1=x_2,\ \dot x_2=u$（位置-速度，对应惯性负载或 $1\,\mathrm{DoF}$ 关节的纯惯量模型），$|u|\le1$，最短时间到原点。由 \cref{thm:oc-bangbang}，$\mathbf B^\top\boldsymbol\lambda=\lambda_2$，故 $u^*=-\operatorname{sgn}[\lambda_2(t)]$。协态解 $\lambda_1=\lambda_{10},\ \lambda_2=\lambda_{20}-\lambda_{10}t$ 是 $t$ 的\emph{一次}函数，故 $\lambda_2(t)$ 至多变号一次——\textbf{$u^*$ 至多切换一次}，取值序列只可能是 $\{-1\},\{+1\},\{-1,+1\},\{+1,-1\}$ 之一。

\paragraph{状态轨线：两族抛物线。}
当 $u^*=+1$：积分得 $x_2=x_{20}+t,\ x_1=x_{10}+x_{20}t+\tfrac12t^2$，消去 $t$ 得轨线 $x_1=\tfrac12x_2^2+C$（开口向右的抛物线族，沿轨线自下而上）。当 $u^*=-1$：$x_1=-\tfrac12x_2^2+C$（开口向左，自上而下）。每族里恰有一条\emph{半支}穿过原点：
\begin{equation}
  \gamma_+:\ x_1=\tfrac12x_2^2\ (x_2\le0),\qquad\gamma_-:\ x_1=-\tfrac12x_2^2\ (x_2\ge0).
  \label{eq:oc-switch-halves}
\end{equation}
把 $\gamma_+,\gamma_-$ 拼成一条过原点的\emph{开关曲线}
\begin{equation}
  \boxed{\gamma:\ h(x_1,x_2)=x_1+\tfrac12x_2|x_2|=0.}
  \label{eq:oc-switch-curve}
\end{equation}

\begin{figure}[ht]
\centering
\begin{tikzpicture}[scale=1.35,>=Stealth]
  \draw[->] (-2.2,0)--(2.2,0) node[right]{$x_2$};
  \draw[->] (0,-1.7)--(0,1.7) node[above]{$x_1$};
  % switching curve gamma: x1 = -1/2 x2|x2|
  \draw[main, very thick, domain=0:1.75, smooth, variable=\s] plot ({\s},{-0.5*\s*\s});
  \draw[main, very thick, domain=-1.75:0, smooth, variable=\s] plot ({\s},{0.5*\s*\s});
  \node[main] at (1.45,-1.05) {$\gamma_+$};
  \node[main] at (-1.45,1.05) {$\gamma_-$};
  \node[main,font=\scriptsize] at (1.95,-0.55){开关曲线 $\gamma$};
  % a sample u=+1 parabola (opens right): x1 = 1/2 x2^2 + C, C=-0.6
  \draw[second, thick, domain=-1.55:1.55, smooth, variable=\s] plot ({\s},{0.5*\s*\s-0.85});
  \node[second,font=\scriptsize] at (1.15,0.55){$u^*=+1$};
  % arrow up on this parabola
  \draw[second,->] (0.4,{0.5*0.16-0.85}) -- (0.55,{0.5*0.30-0.85});
  % a sample u=-1 parabola (opens left): x1 = -1/2 x2^2 + C, C=0.85
  \draw[third, thick, domain=-1.55:1.55, smooth, variable=\s] plot ({\s},{-0.5*\s*\s+0.95});
  \node[third,font=\scriptsize] at (-1.15,-0.55){$u^*=-1$};
  \draw[third,->] (-0.4,{-0.5*0.16+0.95}) -- (-0.55,{-0.5*0.30+0.95});
  % regions
  \node[font=\scriptsize] at (-1.1,-1.25){$R_+$ ($h<0$)};
  \node[font=\scriptsize] at (1.1,1.25){$R_-$ ($h>0$)};
  \fill (0,0) circle (1.2pt);
\end{tikzpicture}
\caption{双积分系统最短时间控制的相平面。开关曲线 $\gamma$（粗）由两条过原点的半抛物线 $\gamma_+,\gamma_-$ 拼成，把 $x_1$--$x_2$ 平面分成 $R_-$（$h>0$）与 $R_+$（$h<0$）。任一初态先沿一族抛物线（细线）运动，撞上 $\gamma$ 后切换控制、沿 $\gamma$ 滑入原点——全程至多切换一次。（仿刘豹图 6.16/6.17 重绘）}
\label{fig:oc-switch-curve}
\end{figure}

\paragraph{最优控制律：状态反馈实现。}
开关曲线把平面分成 $R_-$（$h>0$，曲线右上方）与 $R_+$（$h<0$，左下方）。凡不在 $\gamma$ 上的初态，须先沿当前抛物线运动、撞上 $\gamma$ 后切换一次，再沿 $\gamma$ 滑入原点。综合为状态反馈律：
\begin{equation}
  u^*(t)=\begin{cases}-1,&h(x_1,x_2)>0\quad(\text{即 }(x_1,x_2)\in R_-\cup\gamma_-),\\-\operatorname{sgn}[x_2],&h(x_1,x_2)=0,\\+1,&h(x_1,x_2)<0\quad(\text{即 }(x_1,x_2)\in R_+\cup\gamma_+).\end{cases}
  \label{eq:oc-bangbang-law}
\end{equation}

\begin{figure}[ht]
\centering
\begin{tikzpicture}[
  >=Stealth, node distance=8mm,
  blk/.style={draw, rounded corners, align=center, font=\scriptsize, inner sep=3pt, minimum height=7mm, fill=main!6, draw=main!60!black},
  sum/.style={draw, circle, inner sep=1pt, font=\scriptsize},
  ln/.style={->, thick}]
  \node[sum] (s1) {$+$};
  \node[blk, right=10mm of s1] (rel) {继电器 $R$\\$u^*=-\operatorname{sgn}(z)$};
  \node[blk, right=10mm of rel] (plant) {受控对象\\$\dot x_1=x_2,\dot x_2=u$};
  \node[right=12mm of plant] (out) {};
  \node[blk, below=9mm of rel] (nl) {非线性元件 $N$\\$\tfrac12 x_2|x_2|$};
  \draw[ln] (s1)--node[above,font=\scriptsize]{$z$} (rel);
  \draw[ln] (rel)--node[above,font=\scriptsize]{$u^*$} (plant);
  \draw[ln] (plant.east)--++(1.0,0) node[right,font=\scriptsize]{$\mathbf x$};
  % feedback x1 to summing
  \draw[ln] (plant.east)++(0.5,0) |- (s1) ;
  \node[font=\scriptsize] at ($(plant.east)+(0.5,-0.35)$){$x_1$};
  % x2 to nonlinearity then to sum
  \draw[ln] ($(plant.south)$) |- (nl.east);
  \node[font=\scriptsize] at ($(plant.south)+(0.2,-0.25)$){$x_2$};
  \draw[ln] (nl.west) -| (s1);
\end{tikzpicture}
\caption{Bang--Bang 控制律 \cref{eq:oc-bangbang-law} 的工程实现。状态 $x_2$ 经非线性元件 $N$ 产生 $\tfrac12x_2|x_2|$，与 $x_1$ 相加形成切换信号 $z=h(x_1,x_2)$，再反相经继电器 $R$ 输出 $u^*=\pm1$。继电器的惯性恰好消除其在零输入处工作的不确定性。（仿刘豹图 6.18 重绘）}
\label{fig:oc-bangbang-impl}
\end{figure}

\paragraph{最优转移时间：按控制序列分段求和。}
开关曲线告诉我们\emph{怎么走}，还差一问：从初态到原点究竟\emph{要多久}？这个最短时间本身就是性能指标 $t_f^*=J^*$。计算它的基本方法是\textbf{把最优轨线按控制序列分成若干段、逐段求所需时间、再求和}\cite{liubao2006modern}。双积分情形最简：除已落在开关曲线上的初态外，任一初态的最优轨线恰由\emph{两段}构成——先沿一族抛物线运动到与开关曲线 $\gamma$ 的交点（第一段 $t_a$），再沿 $\gamma$ 滑入原点（第二段 $t_b$），故 $t^*=t_a+t_b$。

以 $R_-$（$h>0$，首段取 $u^*=-1$）为例逐段算。第一段沿 $u=-1$ 抛物线 $x_1=-\tfrac12x_2^2+C_1$（由初态定 $C_1=x_1+\tfrac12x_2^2$）运动，与 $\gamma_+:x_1=\tfrac12x_2^2\,(x_2\le0)$ 相交于 $x_2=\xi=-\sqrt{x_1+\tfrac12x_2^2}$；由 $\dot x_2=u=-1$ 知 $x_2(t)=x_2-t$，故到达交点用时 $t_a=x_2-\xi=x_2+\sqrt{x_1+\tfrac12x_2^2}$。第二段沿 $\gamma_+$（此处 $u=+1$）使 $x_2$ 从 $\xi$ 升回 $0$，用时 $t_b=-\xi=\sqrt{x_1+\tfrac12x_2^2}$。两段相加：
\begin{equation}
  t^*=t_a+t_b=x_2+2\sqrt{x_1+\tfrac12x_2^2}=x_2+\sqrt{4x_1+2x_2^2}.
  \label{eq:oc-bangbang-time-derive}
\end{equation}
$R_+$（首段 $u^*=+1$）与初态已在 $\gamma$ 上（单段直接滑入，用时 $|x_2|$）的情形同理推导，合写为从初态沿最优轨线到达原点的最短时间闭式
\begin{equation}
  t^*(x_1,x_2)=\begin{cases}x_2+\sqrt{4x_1+2x_2^2},&(x_1,x_2)\in R_-,\\|x_2|,&(x_1,x_2)\in\gamma,\\-x_2+\sqrt{-4x_1+2x_2^2},&(x_1,x_2)\in R_+.\end{cases}
  \label{eq:oc-bangbang-time}
\end{equation}

\begin{example}[双积分最短时间：两段相加]\label{ex:oc-bangbang-time}
初态 $\mathbf x_0=(1,0)^\top$。先判区域：$h(1,0)=1+\tfrac12\cdot0=1>0$，属 $R_-$，故首段 $u^*=-1$。\emph{第一段}沿 $x_1=-\tfrac12x_2^2+1$ 运动（$C_1=1$），交 $\gamma_+$ 于 $-\tfrac12\xi^2+1=\tfrac12\xi^2\Rightarrow\xi=-1$，即交点 $(\tfrac12,-1)$；由 $x_2(t)=-t$ 得 $t_a=1$。\emph{第二段}沿 $\gamma_+$（$u^*=+1$）从 $x_2=-1$ 升到 $0$，得 $t_b=1$。总转移时间 $t^*=t_a+t_b=2$，与闭式 \cref{eq:oc-bangbang-time} 的 $x_2+\sqrt{4x_1+2x_2^2}=0+\sqrt4=2$ 完全吻合。对应控制序列为 $\{u^*=-1\ (0\le t<1),\ u^*=+1\ (1\le t<2)\}$，恰一次切换——印证 \cref{sec:oc-double-int-bb} 的“至多切换一次”。
\end{example}

\begin{insight}
\textbf{相平面 + 开关曲线 = 最优反馈的“地图”。}\; 极小值原理给的是\emph{开环}解（$u^*(t)$ 关于时间的函数，需先知 $\boldsymbol\lambda_0$）。但 \cref{fig:oc-switch-curve} 的几何让我们把它翻译成\emph{闭环}状态反馈 \cref{eq:oc-bangbang-law}——只看当前状态落在开关曲线哪一侧就知道该 $+1$ 还是 $-1$。这种“从开环最优解读出反馈律”的思路，正是从最优控制走向可实现控制器的关键一跃，也与 \cref{sec:oc-lqr} 中 LQR“把开环最优解写成 $\mathbf u=-\mathbf K\mathbf x$”一脉相承。机器人学中近似时间最优的“梯形/S 曲线速度规划”即此原理的工程化身。
\end{insight}

\begin{pitfall}[把 Bang--Bang 当“越猛越好”]
Bang--Bang 之所以最快，是因为\emph{时间}是唯一代价、控制能量不计。一旦代价里加入控制能量项（如 LQR 的 $\mathbf u^\top\mathbf R\mathbf u$），最优控制就不再贴边界、而成为光滑的 $-\mathbf K\mathbf x$（\cref{sec:oc-lqr}）。真实机器人几乎从不用纯 Bang--Bang：硬切换激发结构振动、磨损执行器、放大模型误差。时间最优是\emph{理论上界}，工程上常以软化版本（带能量惩罚或速率约束）逼近。
\end{pitfall}

% ════════════════════════════════════════════════════════════════
\section{动态规划与 HJB 方程：全局视角}
\label{sec:oc-dp}

\paragraph{动机：化整为零、逆向递推。}
变分法与极小值原理给的是“沿一条最优轨线的局部必要条件”。动态规划（Bellman，1950 年代）换一个全局视角：把多段决策问题拆成一串单段子问题，从终点\emph{逆向}逐段求解\cite{bellman1957dynamic,liubao2006modern}。它与极小值原理一样能处理控制受界，且天生给出\emph{反馈}形式的解。

\subsection{最优性原理与离散动态规划}
\label{sec:oc-dp-discrete}

\begin{theorem}[Bellman 最优性原理]\label{thm:oc-bellman-principle}
多段决策过程的最优策略具有这样的性质：\emph{不论初始状态与初始决策如何，其余各段决策对于由初始决策所形成的状态而言，也必构成一个最优策略}。等价地：最优轨线的任一后段（从中间任一点到终点）本身也是最优的。
\end{theorem}

\begin{example}[最优路线：穷举 vs 动态规划]\label{ex:oc-route}
汽车从 A 到 B 须穿三条河，每河两桥 P/Q 可选，各段耗时已标注（\cref{fig:oc-route}）。穷举 $8$ 条路线需 $24$ 次加法 + $7$ 次比较，且 $n$ 段时加法数爆增为 $(n-1)2^{n-1}$。动态规划从终点 B 逆推：先比较 $P_3,Q_3$ 到 B 的耗时（得 $Q_3B=2$），再算 $P_2,Q_2$ 到 B 的最优值并比较，逐段前推，仅需 $10$ 次加法 + $6$ 次比较，$n$ 段时为 $4(n-2)+2$ 次加法。$n=10$ 时穷举需 $4608$ 次加法，动态规划仅 $34$ 次。
\end{example}

\begin{figure}[ht]
\centering
\begin{tikzpicture}[
  >=Stealth, node distance=14mm,
  st/.style={draw, circle, font=\scriptsize, inner sep=1.5pt, fill=main!8, draw=main!60!black, minimum size=6mm},
  ed/.style={font=\scriptsize}]
  \node[st] (A) {A};
  \node[st, above right=6mm and 16mm of A] (P1) {$P_1$};
  \node[st, below right=6mm and 16mm of A] (Q1) {$Q_1$};
  \node[st, above right=6mm and 16mm of Q1] (P2) {$P_2$};
  \node[st, below right=6mm and 16mm of Q1] (Q2) {$Q_2$};
  \node[st, above right=6mm and 16mm of Q2] (P3) {$P_3$};
  \node[st, below right=6mm and 16mm of Q2] (Q3) {$Q_3$};
  \node[st, below right=6mm and 16mm of P3] (B) {B};
  \draw (A)--node[ed,above]{3} (P1);
  \draw (A)--node[ed,below]{1} (Q1);
  \draw (P1)--node[ed,above]{6} (P2);
  \draw (P1)--node[ed,left,pos=0.3]{8} (Q2);
  \draw (Q1)--node[ed,left,pos=0.3]{7} (P2);
  \draw (Q1)--node[ed,below]{5} (Q2);
  \draw (P2)--node[ed,above]{6} (P3);
  \draw (P2)--node[ed,left,pos=0.3]{4} (Q3);
  \draw (Q2)--node[ed,left,pos=0.3]{7} (P3);
  \draw (Q2)--node[ed,below]{5} (Q3);
  \draw (P3)--node[ed,above]{4} (B);
  \draw (Q3)--node[ed,below]{2} (B);
  % highlight optimal A-Q1-P2-Q3-B
  \draw[main, very thick] (A)--(Q1)--(P2)--(Q3)--(B);
\end{tikzpicture}
\caption{最优路线决策（\cref{ex:oc-route}）。粗线为动态规划逆推得到的最优路线 $A\,Q_1\,P_2\,Q_3\,B$（总耗时 $12$）。从 B 逐站向前算“各站至终点的最优值”，每步只需一次局部比较，计算量从穷举的指数级降为线性级。（仿刘豹图 6.20/6.21 重绘）}
\label{fig:oc-route}
\end{figure}

\paragraph{离散 Bellman 递推方程。}
对离散系统 $\mathbf x_{k+1}=\mathbf f[\mathbf x_k,\mathbf u_k]$、性能 $J=\Phi[\mathbf x_N]+\sum_{k=0}^{N-1}L[\mathbf x_k,\mathbf u_k]$，记 $J_{N-k}^*[\mathbf x_k]$ 为从 $\mathbf x_k$ 起的后 $N-k$ 段最优代价。由最优性原理得\emph{Bellman 方程}
\begin{equation}
  \boxed{J_{N-k}^*[\mathbf x_k]=\min_{\mathbf u_k}\big\{L[\mathbf x_k,\mathbf u_k]+J_{N-(k+1)}^*[\mathbf x_{k+1}]\big\},\quad\mathbf x_{k+1}=\mathbf f[\mathbf x_k,\mathbf u_k],}
  \label{eq:oc-bellman-discrete}
\end{equation}
边界 $J_0^*[\mathbf x_N]=\Phi[\mathbf x_N]$。从 $k=N-1$ 逆推到 $k=0$，逐段解 $N$ 个单段极小问题，即得最优控制序列 $\{\mathbf u_k^*\}$。

\begin{insight}
\cref{eq:oc-bellman-discrete} 与导论 \cref{eq:ctrl-bellman} 是\emph{同一个}方程——这就是离散 LQR 动态规划推导（导论 \cref{thm:ctrl-lqr-dp}）的出发点。导论在那里令 $L=\mathbf x^\top\mathbf Q\mathbf x+\mathbf u^\top\mathbf R\mathbf u$，归纳证明 $J_t^*=\mathbf x^\top\mathbf P_t\mathbf x$ 二次、$\mathbf u_t^*=-\mathbf K_t\mathbf x$ 线性，得离散 Riccati 递推。本章不重复那段归纳（\cref{sec:oc-lqr-discrete-preview} 仅以标量 DP 例题点醒“反馈增益是状态系数”这一直觉），而把笔墨留给\emph{连续}极限——HJB。
\end{insight}

\begin{example}[一阶离散 LQ：DP 给出状态反馈]\label{ex:oc-dp-scalar}
一阶系统 $x(k+1)=x(k)+u(k),\ x(0)=x_0$，$\min J=\tfrac12cx^2(N)+\tfrac12\sum_{k=0}^{N-1}u^2(k)$（$c>0$）。取 $N=2$。\emph{逆推第一步}（末段 $x(1)\to x(2)$）：$J_1[x(1)]=\tfrac12u^2(1)+\tfrac12c[x(1)+u(1)]^2$，令 $\partial J_1/\partial u(1)=0$ 得 $u^*(1)=-\tfrac{c}{1+c}x(1)$，$J_1^*[x(1)]=\tfrac{c}{2(1+c)}x^2(1)$。\emph{逆推第二步}（$x(0)\to x(1)$）：$J_2[x(0)]=\tfrac12u^2(0)+J_1^*[x(0)+u(0)]$，令导数为零得 $u^*(0)=-\tfrac{c}{1+2c}x(0)$，$J_2^*[x(0)]=\tfrac{c}{2(1+2c)}x^2(0)$。
\textbf{两点观察}：(i) 每段最优控制都是\emph{当段状态的线性函数} $u^*(k)=-f_kx(k)$——离散 LQR 的反馈本质；(ii) 最优 cost-to-go 都是状态的\emph{二次型}。这正是导论 \cref{thm:ctrl-lqr-dp} 在标量、一般 $(a,b,q,r)$ 下的特例。
\end{example}

\subsection{离散最优控制的第二条路：拉格朗日乘子与离散正则方程}
\label{sec:oc-discrete-canonical}

\paragraph{动机：把离散问题当成一道放大了 $N$ 倍的静态最优化。}
Bellman 递推 \cref{eq:oc-bellman-discrete} 从\emph{值函数}的全局视角解离散问题。但对\emph{同一个}离散问题还有一条与之并行、却更贴近\emph{变分/协态}的局部路线：既然把 $[t_0,t_f]$ 分成 $N$ 段后，待优化的 $\{\mathbf x(k)\},\{\mathbf u(k)\}$ 只是\emph{有限}多个变量，整个问题就退化为一道“变量数放大 $N$ 倍”的等式约束静态最优化——正可照搬 \cref{sec:oc-static} 的 Lagrange 乘子法，得到连续 Hamilton 正则方程 \cref{thm:oc-canonical} 的\emph{离散孪生}\cite{liubao2006modern}。把这条路与上面的 Bellman 路并置，离散世界就和连续世界一样有了“值函数（DP）”与“协态（变分/极小值原理）”两套等价语言；且这条路不要求值函数处处可微（那是 \cref{sec:oc-hjb} HJB 的软肋），直接给出离散极小值原理，是数值最优控制（间接法/打靶法）的算法地基。

\paragraph{问题、约束与增广指标。}
设离散系统与综合型性能指标
\begin{align}
  \mathbf x(k+1)&=\mathbf f[\mathbf x(k),\mathbf u(k),k],\quad \mathbf x(0)=\mathbf x_0,\quad k=0,\dots,N-1,\label{eq:oc-disc-sys}\\
  J&=\Phi[\mathbf x(N)]+\sum_{k=0}^{N-1}L[\mathbf x(k),\mathbf u(k),k].\label{eq:oc-disc-cost}
\end{align}
把状态方程视为每段一条的等式约束 $\mathbf f[\mathbf x(k),\mathbf u(k),k]-\mathbf x(k+1)=\mathbf0$，对每段引入 $n$ 维乘子 $\boldsymbol\lambda(k+1)$（$k=0,\dots,N-1$，即乘子序列 $\boldsymbol\lambda(1),\dots,\boldsymbol\lambda(N)$），构造增广指标
\begin{equation}
  V=\Phi[\mathbf x(N)]+\sum_{k=0}^{N-1}\Big\{L[\mathbf x(k),\mathbf u(k),k]+\boldsymbol\lambda^\top(k+1)\big(\mathbf f[\mathbf x(k),\mathbf u(k),k]-\mathbf x(k+1)\big)\Big\}.
  \label{eq:oc-disc-augmented}
\end{equation}
这就把“在状态方程约束下求 $J$ 极小”化成了“对 $V$ 求无约束极小”。定义\emph{离散 Hamilton 函数}
\begin{equation}
  \boxed{H_k=L[\mathbf x(k),\mathbf u(k),k]+\boldsymbol\lambda^\top(k+1)\,\mathbf f[\mathbf x(k),\mathbf u(k),k],}
  \label{eq:oc-disc-hamiltonian}
\end{equation}
\textbf{它用的是\emph{下一步}的协态 $\boldsymbol\lambda(k+1)$}——这是离散框架区别于连续 $H=L+\boldsymbol\lambda^\top\mathbf f$ 的关键记号细节，下面的“逆向递推”正源于此。

\begin{derivation}[从 $\delta V=0$ 挤出离散正则方程：分部求和]
代入 \cref{eq:oc-disc-hamiltonian}，$V=\Phi[\mathbf x(N)]+\sum_{k=0}^{N-1}\{H_k-\boldsymbol\lambda^\top(k+1)\mathbf x(k+1)\}$。对末项作\emph{分部求和}（连续分部积分的离散版，把指标 $k+1$ 平移回 $k$）：$\sum_{k=0}^{N-1}\boldsymbol\lambda^\top(k+1)\mathbf x(k+1)=\sum_{k=1}^{N-1}\boldsymbol\lambda^\top(k)\mathbf x(k)+\boldsymbol\lambda^\top(N)\mathbf x(N)$，再把 $\sum_{k=0}^{N-1}H_k=H_0+\sum_{k=1}^{N-1}H_k$ 拆出首项，得
\begin{equation}
  V=\Phi[\mathbf x(N)]-\boldsymbol\lambda^\top(N)\mathbf x(N)+H_0+\sum_{k=1}^{N-1}\big[H_k-\boldsymbol\lambda^\top(k)\mathbf x(k)\big].
  \label{eq:oc-disc-Vrearr}
\end{equation}
求 $V$ 的一阶增量、略去高阶项（始端 $\mathbf x(0)=\mathbf x_0$ 固定故 $\Delta\mathbf x(0)=\mathbf0$，其 $\partial H_0/\partial\mathbf x(0)$ 项自动消失）：
\begin{equation}
  \Delta V=\Delta\mathbf x^\top(N)\Big[\frac{\partial\Phi}{\partial\mathbf x(N)}-\boldsymbol\lambda(N)\Big]+\sum_{k=1}^{N-1}\Delta\mathbf x^\top(k)\Big[\frac{\partial H_k}{\partial\mathbf x(k)}-\boldsymbol\lambda(k)\Big]+\sum_{k=0}^{N-1}\Delta\mathbf u^\top(k)\,\frac{\partial H_k}{\partial\mathbf u(k)}.
  \label{eq:oc-disc-variation}
\end{equation}
极小要求 $\Delta V=0$ 对一切独立变分 $\Delta\mathbf x(k)\,(k=1,\dots,N)$、$\Delta\mathbf u(k)\,(k=0,\dots,N-1)$ 成立——这正是 \cref{thm:oc-variation-zero} 的离散版（“变分为零”在有限维退回“偏导为零”）。遂令各方括号与系数分别为零，即得正则方程。
\end{derivation}

\begin{theorem}[离散时间正则方程]\label{thm:oc-disc-canonical}
离散最优控制问题 \cref{eq:oc-disc-sys,eq:oc-disc-cost}（控制\emph{无界}）的最优解 $\{\mathbf x^*(k),\mathbf u^*(k),\boldsymbol\lambda^*(k)\}$ 满足
\begin{align}
  \mathbf x(k+1)&=\frac{\partial H_k}{\partial\boldsymbol\lambda(k+1)}=\mathbf f[\mathbf x(k),\mathbf u(k),k] &&\text{(状态方程，\emph{正}向递推)},\label{eq:oc-disc-state}\\
  \boldsymbol\lambda(k)&=\frac{\partial H_k}{\partial\mathbf x(k)}=\frac{\partial L}{\partial\mathbf x(k)}+\Big[\frac{\partial\mathbf f}{\partial\mathbf x(k)}\Big]^\top\boldsymbol\lambda(k+1) &&\text{(协态方程，\emph{逆}向递推)},\label{eq:oc-disc-costate}\\
  \mathbf0&=\frac{\partial H_k}{\partial\mathbf u(k)}=\frac{\partial L}{\partial\mathbf u(k)}+\Big[\frac{\partial\mathbf f}{\partial\mathbf u(k)}\Big]^\top\boldsymbol\lambda(k+1) &&\text{(控制方程/平稳条件)},\label{eq:oc-disc-control}
\end{align}
配两点边界 $\mathbf x(0)=\mathbf x_0$（始端）与 $\boldsymbol\lambda(N)=\partial\Phi/\partial\mathbf x(N)$（终端横截条件）。状态由 $k=0$ 正推、协态由 $k=N$ 逆推，二者在 $[0,N]$ 上耦合，构成离散\emph{两点边值问题}。
\end{theorem}

\noindent 规模核对（呼应“放大 $N$ 倍的静态优化”这一动机）：\cref{eq:oc-disc-costate,eq:oc-disc-control,eq:oc-disc-state} 与两条边界依次给出 $n(N-1)$、$rN$、$nN$、$n$、$n$ 个标量方程，共 $(2n+r)N+n$ 个；而未知量 $\mathbf x,\boldsymbol\lambda,\mathbf u$（除去已知的 $\mathbf x(0)$ 这 $n$ 个）恰为 $(2n+r)N$ 个——方程比未知数多出的 $n$ 个，正好被 $n$ 维终端条件 $\boldsymbol\lambda(N)=\partial\Phi/\partial\mathbf x(N)$ 吸收，问题适定。

\begin{insight}
\textbf{离散正则方程是连续 Hamilton 正则方程 \cref{thm:oc-canonical} 的逐行孪生。}\;
两套必要条件一一对应，差别只在“微分 $\leftrightarrow$ 差分递推”与协态用的是 $\boldsymbol\lambda(k+1)$：
\begin{center}
\small
\begin{tabular}{lll}
\toprule
 & 连续（\cref{thm:oc-canonical}）& 离散（\cref{thm:oc-disc-canonical}）\\
\midrule
Hamilton & $H=L+\boldsymbol\lambda^\top\mathbf f$ & $H_k=L+\boldsymbol\lambda^\top(k{+}1)\mathbf f$ \\
状态 & $\dot{\mathbf x}=+\partial H/\partial\boldsymbol\lambda$ & $\mathbf x(k{+}1)=\partial H_k/\partial\boldsymbol\lambda(k{+}1)$ \\
协态 & $\dot{\boldsymbol\lambda}=-\partial H/\partial\mathbf x$ & $\boldsymbol\lambda(k)=\partial H_k/\partial\mathbf x(k)$（逆推）\\
平稳 & $\partial H/\partial\mathbf u=\mathbf0$ & $\partial H_k/\partial\mathbf u(k)=\mathbf0$ \\
终端 & $\boldsymbol\lambda(t_f)=\partial\Phi/\partial\mathbf x(t_f)$ & $\boldsymbol\lambda(N)=\partial\Phi/\partial\mathbf x(N)$ \\
\bottomrule
\end{tabular}
\end{center}
取连续极限可验证二者自洽：令 $\mathbf f=\mathbf x+\Delta t\,\mathbf a(\mathbf x,\mathbf u)$（$\dot{\mathbf x}=\mathbf a$）、$L=\Delta t\,L_c$，则 $\partial\mathbf f/\partial\mathbf x=\mathbf I+\Delta t\,\partial\mathbf a/\partial\mathbf x$，协态递推 \cref{eq:oc-disc-costate} 化为 $\big(\boldsymbol\lambda(k)-\boldsymbol\lambda(k+1)\big)/\Delta t=\partial L_c/\partial\mathbf x+(\partial\mathbf a/\partial\mathbf x)^\top\boldsymbol\lambda$，令 $\Delta t\to0$ 即 $\dot{\boldsymbol\lambda}=-\partial H_c/\partial\mathbf x$——回到连续协态方程。这正是“先离散化、求解、再令 $\Delta t\to0$”这条处理思路（刘豹 §6.5）的合法性根据。
\end{insight}

\paragraph{用于线性二次型：殊途同归于离散 Riccati。}
把 \cref{thm:oc-disc-canonical} 用在线性系统 $\mathbf x(k+1)=\mathbf G\mathbf x(k)+\mathbf H\mathbf u(k)$、二次代价 $J=\tfrac12\sum[\mathbf x^\top\mathbf Q\mathbf x+\mathbf u^\top\mathbf R\mathbf u]+\tfrac12\mathbf x^\top(N)\mathbf Q_f\mathbf x(N)$ 上：$H_k=\tfrac12[\mathbf x^\top\mathbf Q\mathbf x+\mathbf u^\top\mathbf R\mathbf u]+\boldsymbol\lambda^\top(k+1)[\mathbf G\mathbf x+\mathbf H\mathbf u]$，协态方程给 $\boldsymbol\lambda(k)=\mathbf Q\mathbf x(k)+\mathbf G^\top\boldsymbol\lambda(k+1)$，控制方程给 $\mathbf u^*(k)=-\mathbf R^{-1}\mathbf H^\top\boldsymbol\lambda(k+1)$。再设扫描 ansatz $\boldsymbol\lambda(k)=\mathbf P(k)\mathbf x(k)$ 并消元，便逆推出\emph{离散 Riccati 递推}与反馈增益——这与导论 \cref{thm:ctrl-lqr-dp} 用\emph{动态规划}得到的离散 Riccati \cref{eq:ctrl-dlqr-P}、DARE \cref{eq:ctrl-dare} \emph{完全一致}，只是出发点从“值函数二次型”换成了“协态扫描”。故本章不重推该归纳（见导论），此处只点明：连续侧的“协态 $\leftrightarrow$ DP”二象（\cref{sec:oc-hjb} 的 $\boldsymbol\lambda=\partial J^*/\partial\mathbf x$）在离散侧同样成立。

\begin{pitfall}[离散“极小值原理”不能无条件照搬连续版]
\cref{thm:oc-disc-canonical} 给的是控制\emph{无界}时的平稳形式 $\partial H_k/\partial\mathbf u(k)=\mathbf0$。当控制受界 $\mathbf u(k)\in U$ 时，人们自然想照搬连续极小值原理 \cref{eq:oc-pmp-min}、写成 $\mathbf u^*(k)=\arg\min_{\mathbf u\in U}H_k$。\textbf{但这一“离散极小值原理”并非无条件成立}——连续 PMP 的针状变分证明依赖了“控制集生成的瞬时可达方向锥是凸的”，这在连续时间因时间可无限细分而自动满足，离散时间却\emph{未必}；缺了附加凸性（如可达集的方向凸性）时，$\arg\min H_k$ 只是更弱意义下的必要条件，可能给出非最优解。故离散问题受界时须谨慎，严格陈述与反例属高级内容，留待 \cref{ch:opt-variational-pmp}。本章在此只给无界情形的平稳形式（与刘豹 §6.4.1 一致）。
\end{pitfall}

\subsection{连续动态规划：Hamilton--Jacobi--Bellman 方程}
\label{sec:oc-hjb}

\paragraph{从离散递推到偏微分方程。}
令段长 $\Delta t\to0$，离散 Bellman 递推过渡到连续形式。设 $J^*[\mathbf x,t]=\min_{\mathbf u\in U}\{\int_t^{t_f}L\,\mathrm dt+\Phi[\mathbf x(t_f)]\}$，取 $t'=t+\Delta t$。由最优性原理
\begin{equation}
  J^*[\mathbf x,t]=\min_{\mathbf u\in U}\big\{L[\mathbf x,\mathbf u,t]\,\Delta t+J^*[\mathbf x(t+\Delta t),t+\Delta t]\big\}.
  \label{eq:oc-hjb-step}
\end{equation}
把 $J^*[\mathbf x+\Delta\mathbf x,t+\Delta t]$ 在 $[\mathbf x,t]$ 处 Taylor 展开一阶（$\Delta\mathbf x=\mathbf f\,\Delta t$）：$J^*[\mathbf x,t]+(\partial J^*/\partial\mathbf x)^\top\mathbf f\,\Delta t+(\partial J^*/\partial t)\,\Delta t$，代入并约去 $J^*[\mathbf x,t]$、令 $\Delta t\to0$，得

\begin{theorem}[Hamilton--Jacobi--Bellman 方程]\label{thm:oc-hjb}
连续动态规划基本方程（HJB）为
\begin{equation}
  \boxed{-\frac{\partial J^*[\mathbf x,t]}{\partial t}=\min_{\mathbf u\in U}\Big\{L[\mathbf x,\mathbf u,t]+\Big[\frac{\partial J^*[\mathbf x,t]}{\partial\mathbf x}\Big]^\top\mathbf f[\mathbf x,\mathbf u,t]\Big\},}
  \label{eq:oc-hjb}
\end{equation}
边界条件 $J^*[\mathbf x(t_f),t_f]=\Phi[\mathbf x(t_f),t_f]$。若令 $\boldsymbol\lambda=\partial J^*/\partial\mathbf x$，则花括号内恰为 Hamilton 函数 $H[\mathbf x,\mathbf u,\boldsymbol\lambda,t]$，故 HJB 可写成 $-\partial J^*/\partial t=\min_{\mathbf u\in U}H[\mathbf x,\mathbf u,\boldsymbol\lambda,t]$。
\end{theorem}

\paragraph{HJB 与极小值原理互为表里。}
HJB 的 $\min_{\mathbf u\in U}H$ 与极小值原理 \cref{eq:oc-pmp-min} 完全一致——\textbf{HJB 实际上是极小值原理的另一种形式}。更进一步，对 HJB 关于 $\mathbf x$ 求偏导，并用 $\boldsymbol\lambda=\partial J^*/\partial\mathbf x$、$\frac{\mathrm d}{\mathrm dt}\partial_\mathbf x J^*=\partial_{\mathbf x t}^2J^*+\partial_{\mathbf x\mathbf x}^2J^*\dot{\mathbf x}$，可\emph{推导出}协态方程 $\dot{\boldsymbol\lambda}=-\partial H/\partial\mathbf x$ 与横截条件——这等于用动态规划\emph{间接证明}了极小值原理\cite{liubao2006modern}。三者（变分/协态、极小值原理、动态规划）至此完全打通：

\begin{insight}
\textbf{协态 $=$ 值函数梯度。}\; 关系 $\boldsymbol\lambda(t)=\dfrac{\partial J^*[\mathbf x,t]}{\partial\mathbf x}$ 是连接三套语言的“罗塞塔石碑”：
\begin{center}
\small
\begin{tabular}{lll}
\toprule
框架 & 核心对象 & 核心方程 \\
\midrule
变分 / 协态（PMP 局部）& 协态 $\boldsymbol\lambda(t)$ & 正则方程 + $\min_{\mathbf u}H$ \\
动态规划（全局）& 值函数 $J^*[\mathbf x,t]$ & HJB 偏微分方程 \\
桥梁 & $\boldsymbol\lambda=\partial J^*/\partial\mathbf x$ & 二者等价 \\
\bottomrule
\end{tabular}
\end{center}
代价是：HJB 是\emph{偏}微分方程，要求 $J^*$ 处处可微——而许多实际问题 $J^*$ 不光滑（如 Bang--Bang 的开关曲线处），这正是 DP 的使用瓶颈，也是 A 章 PMP（针状变分/奇异弧）与 B 章 HJB（黏性解）要严格处理的地方。
\end{insight}

\begin{example}[连续 HJB 求最优反馈]\label{ex:oc-hjb-example}
系统 $\dot{\mathbf x}=\begin{psmallmatrix}0&1\\0&0\end{psmallmatrix}\mathbf x+\begin{psmallmatrix}0\\1\end{psmallmatrix}u$，初态 $\mathbf x(0)=(1,0)^\top$\footnote{$\mathbf x(0)=(1,0)^\top$（刘豹勘误 M2-e：教材 OCR 把列向量两分量颠倒成 $(0,1)^\top$；由最终 $J^*=2x_1^2+2x_1x_2+x_2^2$ 在 $(1,0)$ 取值 $2$、且解 $\mathbf x^*,u^*,J^*$ 用 sympy 复核均只在 $(1,0)^\top$ 成立可证）。}，$J=\int_0^\infty(2x_1^2+\tfrac12u^2)\,\mathrm dt$，$u$ 无界。Hamilton 函数 $H=2x_1^2+\tfrac12u^2+\partial_{x_1}J^*\,x_2+\partial_{x_2}J^*\,u$。由 $\partial H/\partial u=0$ 得 $u^*=-\partial J^*/\partial x_2$。$J^*$ 不显含 $t$（无限时域），代入 HJB 得偏微分方程 $2x_1^2+\partial_{x_1}J^*\,x_2-\tfrac12(\partial_{x_2}J^*)^2=0$。设二次型解 $J^*=a_1x_1^2+2a_2x_1x_2+a_3x_2^2$，代入匹配系数得 $a_2=1,a_3=1,a_1=2$，故
\begin{equation}
  J^*[\mathbf x]=2x_1^2+2x_1x_2+x_2^2,\qquad u^*=-\partial J^*/\partial x_2=-2(x_1+x_2).
  \label{eq:oc-hjb-feedback}
\end{equation}
最优控制是\emph{状态线性反馈}——HJB 在“线性系统 + 二次代价”下自动回到 LQR。验证：闭环 $\dot{\mathbf x}^*=\begin{psmallmatrix}0&1\\-2&-2\end{psmallmatrix}\mathbf x^*$，解得 $\mathbf x^*(t)=(\mathrm e^{-t}(\cos t+\sin t),\,-2\mathrm e^{-t}\sin t)^\top$，$J^*=\int_0^\infty4\mathrm e^{-2t}\,\mathrm dt=2$，与 $J^*[\mathbf x(0)]=2x_1^2|_{(1,0)}=2$ 吻合。
\end{example}

\noindent 当 $H$ 对 $u$ 非二次（如代价含 $|u|$ 项或控制受界），HJB 退化为非线性偏微分方程，一般只能数值求解——例如 $\ddot y+4\dot y+y=u,\ |u|\le1$ 的最短时间问题，HJB 给 $u^*=-\operatorname{sgn}(\partial J^*/\partial x_2)$，但 $J^*$ 须靠数值（电子计算机）求得。

% ════════════════════════════════════════════════════════════════
\section{线性二次型最优控制（LQR）：物理意义、调节器与跟踪器}
\label{sec:oc-lqr}

\paragraph{定位：导论已推透，本章补直觉与跟踪器。}
\cref{ch:control} 的 \cref{sec:ctrl-lqr} 已用动态规划与 HJB 两条路把 LQR 推导得淋漓尽致：离散 Riccati 递推（\cref{eq:ctrl-dlqr-P}）、DARE（\cref{eq:ctrl-dare}）、CARE（\cref{eq:ctrl-care}）、Hamilton 矩阵解法（\cref{eq:ctrl-hamilton}）、闭环稳定性（值函数即 Lyapunov，\cref{thm:ctrl-lqr-lyap}）、乃至 Kalman 反例（\cref{ex:ctrl-kalman}）。\textbf{本节不重复这些推导}，而是补两层导论留白：(1) 二次型代价各项的\emph{物理意义逐项解析}——告诉你权重该怎么选；(2) 导论未涉的\emph{有限时域 Riccati 微分方程}的工程含义、\emph{输出调节器}与\emph{跟踪器}。严格的存在唯一性与博弈型 Riccati 见 \cref{ch:lqr-lqg-riccati}。

\subsection{二次型代价各项的物理意义}
\label{sec:oc-quadratic-meaning}

\paragraph{为什么是二次型？}
若系统线性、性能泛函是状态/控制的二次型积分，则最优控制问题称\emph{线性二次型}（LQ）。它的解最简单——是状态的\emph{线性}函数，可由状态反馈实现——故工程上应用极广\cite{liubao2006modern}。一般形式（连续、有限时域）：
\begin{equation}
  J=\frac12\int_{t_0}^{t_f}\big[\mathbf x^\top\mathbf Q(t)\mathbf x+\mathbf u^\top\mathbf R(t)\mathbf u\big]\mathrm dt+\frac12\mathbf x^\top(t_f)\mathbf Q_f\,\mathbf x(t_f),
  \label{eq:oc-lq-cost}
\end{equation}
其中 $\mathbf Q(t)\succeq\mathbf0$（半正定状态权）、$\mathbf R(t)\succ\mathbf0$（正定控制权）、$\mathbf Q_f\succeq\mathbf0$（半正定终端权），工程上常取对角常阵。\footnote{此处终端项为 $\tfrac12\mathbf x^\top(t_f)\mathbf Q_f\mathbf x(t_f)$（刘豹勘误 M2-e：教材 OCR 把末尾状态因子 $\mathbf x(t_f)$ 误并成 $\mathbf Q_0$ 的宗量 $(t_f)$，但 $\mathbf Q_f$ 是常值终端权、且该项须为标量二次型才有意义）。}

\begin{insight}
\textbf{逐项物理意义（这是导论缺、本章补的直觉层）}\cite{liubao2006modern}：
\begin{itemize}
  \item \textbf{误差代价} $L_x=\tfrac12\mathbf x^\top\mathbf Q\mathbf x$：若 $\mathbf x$ 表误差，则 $L_x$ 是“误差平方”的度量。$\mathbf Q\succeq\mathbf0\Rightarrow$ 只要有误差 $L_x\ge0$；$\mathbf x=\mathbf0\Rightarrow L_x=0$。$\mathbf Q$ 常取对角，对角元 $q_i$ 表“对误差分量 $x_i$ 的重视程度”——\emph{越想压小哪个分量、就把对应 $q_i$ 取越大}。若不同阶段重视程度不同，$q_i$ 可取时变。
  \item \textbf{控制能量代价} $L_u=\tfrac12\mathbf u^\top\mathbf R\mathbf u$：度量控制“功率”。$\mathbf R\succ\mathbf0\Rightarrow$ 只要有控制 $L_u>0$。若 $\mathbf u$ 是电压/电流，则 $\tfrac12\int\mathbf u^\top\mathbf R\mathbf u\,\mathrm dt\propto$ 区间内消耗的\emph{能量}。\emph{$\mathbf R$ 取越大、越“吝啬”用控制}。
  \item \textbf{终端代价} $\tfrac12\mathbf x^\top(t_f)\mathbf Q_f\mathbf x(t_f)$：突出对\emph{终端}误差的要求（如航天器交会要求终态严格一致，则必加此项以压小 $t_f$ 时的误差）。
\end{itemize}
\textbf{一句话}：$J\to\min$ 的实质，是“用不大的控制能量，保持较小的误差”——$\mathbf Q$ 与 $\mathbf R$ 之比设定了“误差 vs 省力”的天平。这层直觉，正是导论与 A--Q 偏重推导而缺的一环。
\end{insight}

\subsection{有限时域状态调节器：Riccati 微分方程}
\label{sec:oc-state-regulator}

\paragraph{从极小值原理到 Riccati 微分方程。}
\emph{状态调节器}的任务：当状态因扰动偏离平衡（取零状态为平衡），用不过多能量把它拉回零点附近。设线性时变系统 $\dot{\mathbf x}=\mathbf A(t)\mathbf x+\mathbf B(t)\mathbf u$，代价 \cref{eq:oc-lq-cost}，$\mathbf u$ 无界。Hamilton 函数 $H=\tfrac12[\mathbf x^\top\mathbf Q\mathbf x+\mathbf u^\top\mathbf R\mathbf u]+\boldsymbol\lambda^\top[\mathbf A\mathbf x+\mathbf B\mathbf u]$。控制方程 $\partial H/\partial\mathbf u=\mathbf R\mathbf u+\mathbf B^\top\boldsymbol\lambda=\mathbf0$ 给
\begin{equation}
  \mathbf u^*=-\mathbf R^{-1}\mathbf B^\top\boldsymbol\lambda,
  \label{eq:oc-lqr-u-costate}
\end{equation}
$\partial^2H/\partial\mathbf u^2=\mathbf R\succ\mathbf0$ 保证极小（必要且充分）。正则方程 $\dot{\boldsymbol\lambda}=-\mathbf Q\mathbf x-\mathbf A^\top\boldsymbol\lambda$，$\dot{\mathbf x}=\mathbf A\mathbf x-\mathbf B\mathbf R^{-1}\mathbf B^\top\boldsymbol\lambda$\footnote{$\dot{\mathbf x}=+\partial H/\partial\boldsymbol\lambda=\mathbf A\mathbf x+\mathbf B\mathbf u$（刘豹勘误 M2-f：教材式 6.310 OCR 误印 $\dot{\mathbf x}=-\partial H/\partial\boldsymbol\lambda$，但 Hamilton 正则方程为 $\dot{\mathbf x}=+\partial H/\partial\boldsymbol\lambda$，与右端 $\mathbf A\mathbf x+\mathbf B\mathbf u$ 才自洽）。}，终端 $\boldsymbol\lambda(t_f)=\mathbf Q_f\mathbf x(t_f)$。

\paragraph{扫描变换：从协态到反馈。}
为把 \cref{eq:oc-lqr-u-costate} 化成状态反馈，设\emph{扫描 ansatz} $\boldsymbol\lambda(t)=\mathbf P(t)\mathbf x(t)$（$\mathbf P$ 为对称正定待定阵）。代入得 $\mathbf u^*=-\mathbf R^{-1}\mathbf B^\top\mathbf P(t)\mathbf x(t)=-\mathbf K(t)\mathbf x(t)$，反馈增益 $\mathbf K(t)=\mathbf R^{-1}\mathbf B^\top\mathbf P(t)$。把 $\boldsymbol\lambda=\mathbf P\mathbf x$ 求导 $\dot{\boldsymbol\lambda}=\dot{\mathbf P}\mathbf x+\mathbf P\dot{\mathbf x}$，与正则方程联立、消去 $\mathbf x$（对所有 $\mathbf x$ 成立），得

\begin{theorem}[连续 Riccati 微分方程（RDE）]\label{thm:oc-rde}
有限时域状态调节器的反馈阵 $\mathbf P(t)$ 满足
\begin{equation}
  \boxed{-\dot{\mathbf P}(t)=\mathbf P(t)\mathbf A(t)+\mathbf A^\top(t)\mathbf P(t)-\mathbf P(t)\mathbf B(t)\mathbf R^{-1}(t)\mathbf B^\top(t)\mathbf P(t)+\mathbf Q(t),\quad\mathbf P(t_f)=\mathbf Q_f,}
  \label{eq:oc-rde}
\end{equation}
从终端 $t_f$ \emph{逆时间}积分；最优控制 $\mathbf u^*=-\mathbf R^{-1}\mathbf B^\top\mathbf P(t)\mathbf x$。\footnote{教材式 6.319 的 OCR 漏了左端导数点 $\dot{\mathbf P}$、并在 $\mathbf B^\top$ 与 $\mathbf P$ 间多出杂散 $\times$（刘豹勘误 M2-f）；正确即上式 RDE，终端条件 $\mathbf P(t_f)=\mathbf Q_f$。}
$\mathbf P(t)$ 对称（由 RDE 转置后与原方程同解、解唯一可证），故只需解 $\tfrac{n(n+1)}{2}$ 个标量一阶方程。
\end{theorem}

\noindent \cref{eq:oc-rde} 与导论 \cref{eq:ctrl-rde} 同形——这是同一个 RDE。导论从 HJB（含时值函数 $J^*=\mathbf x^\top\mathbf P(t)\mathbf x$）导出它，本章从极小值原理（协态扫描）导出它，二者殊途同归。性能最优值有简洁闭式
\begin{equation}
  J^*[\mathbf x(t_0)]=\tfrac12\mathbf x^\top(t_0)\mathbf P(t_0)\mathbf x(t_0),\qquad J^*[\mathbf x(t)]=\tfrac12\mathbf x^\top(t)\mathbf P(t)\mathbf x(t),
  \label{eq:oc-lqr-Jstar}
\end{equation}
即 $\mathbf P(t)$ 正是“从 $t$ 起的最优 cost-to-go 权重”（与导论 cost-to-go 同义，差一个 $\tfrac12$）。

\begin{example}[一阶调节器：Riccati 的“稳态”浮现]\label{ex:oc-lqr-scalar}
一阶系统 $\dot x=ax+u,\ x(0)=x_0$，$J=\tfrac12\int_0^{t_f}[q_1x^2+q_2u^2]\,\mathrm dt+\tfrac12q_0x^2(t_f)$。由 \cref{eq:oc-rde}，标量 Riccati $\dot P=-2aP+\tfrac1{q_2}P^2-q_1,\ P(t_f)=q_0$，最优控制 $u^*=-\tfrac1{q_2}P(t)x$。其闭式含 $\beta=\sqrt{\tfrac{q_1}{q_2}+a^2}$\footnote{$\beta=\sqrt{q_1/q_2+a^2}$（刘豹勘误 M2-f：教材误印 $\dot a^2$，应为 $a^2$，由其 $t_f\to\infty$ 极限 $\sqrt{q_1/q_2+a^2}$ 反证）。}。
\textbf{关键观察}（\cref{fig:oc-riccati}）：从终端 $t_f$ 逆时间看，$P(t)$ 随 $t$ 减小而趋近一个“稳态值”，且\emph{该值与终端条件 $q_0$ 无关}；$t_f$ 越大，$P(t)$ 保持常值的区间越宽。$t_f\to\infty$ 时
\begin{equation}
  \lim_{t_f\to\infty}P(t)=q_2(\beta+a)=aq_2+q_2\sqrt{\tfrac{q_1}{q_2}+a^2}.
  \label{eq:oc-riccati-steady}
\end{equation}
例如 $a=-1,q_1=q_2=1$ 时 $\lim P=-1+\sqrt2=0.414$（一个常数）。
\end{example}

\begin{figure}[ht]
\centering
\begin{tikzpicture}[scale=1.0,>=Stealth]
  \draw[->] (-0.3,0)--(5.3,0) node[right]{$t$};
  \draw[->] (0,-0.3)--(0,2.6) node[above]{$P(t)$};
  % steady value line
  \draw[dashed,gray] (0,1.55)--(5.0,1.55) node[right,font=\scriptsize,gray]{稳态值 (与 $q_0$ 无关)};
  % small tf curve
  \draw[main, thick] (0,1.50) -- (1.4,1.50) .. controls (1.8,1.50) and (2.0,1.30) .. (2.1,0.6);
  \fill[main] (2.1,0.6) circle (1.1pt); \node[main,font=\scriptsize] at (2.1,-0.25){$t_{f1}$};
  \node[main,font=\scriptsize] at (1.4,1.72){小 $t_f$};
  % large tf curve
  \draw[second, thick] (0,1.55) -- (3.6,1.55) .. controls (4.0,1.55) and (4.2,1.35) .. (4.3,0.6);
  \fill[second] (4.3,0.6) circle (1.1pt); \node[second,font=\scriptsize] at (4.3,-0.25){$t_{f2}$};
  \node[second,font=\scriptsize] at (3.3,1.73){大 $t_f$};
\end{tikzpicture}
\caption{有限时域 Riccati 解 $P(t)$ 随终端时刻 $t_f$ 的行为（\cref{ex:oc-lqr-scalar}）。从 $t_f$ 逆时间看，$P(t)$ 迅速趋向一个与终端权 $q_0$ 无关的“稳态值”\cref{eq:oc-riccati-steady}；$t_f$ 越大，常值区间越宽。这解释了为何 $t_f$ 足够大时可用常增益 LQR 近似——即下节的无限时域调节器。（仿刘豹图 6.32 重绘）}
\label{fig:oc-riccati}
\end{figure}

\subsection{无限时域调节器：代数 Riccati 与常增益反馈}
\label{sec:oc-infinite-regulator}

\paragraph{从“稳态”到常增益。}
\cref{fig:oc-riccati} 启发：随 $t_f\to\infty$，$\mathbf P(t)\to$ 常阵 $\mathbf P$，时变最优反馈退化为\emph{定常}反馈。这就是\emph{无限时域状态调节器}\cite{liubao2006modern,anderson2007optimal}：

\begin{theorem}[无限时域 LQR]\label{thm:oc-infinite-lqr}
若线性定常系统 $\dot{\mathbf x}=\mathbf A\mathbf x+\mathbf B\mathbf u$ \emph{能控}，$\mathbf Q\succeq\mathbf0$、$\mathbf R\succ\mathbf0$，$J=\tfrac12\int_{t_0}^{\infty}[\mathbf x^\top\mathbf Q\mathbf x+\mathbf u^\top\mathbf R\mathbf u]\,\mathrm dt$，则最优控制存在且唯一，为常增益状态反馈
\begin{equation}
  \mathbf u^*=-\mathbf R^{-1}\mathbf B^\top\mathbf P\,\mathbf x=-\mathbf K\mathbf x,\qquad\mathbf K=\mathbf R^{-1}\mathbf B^\top\mathbf P,
  \label{eq:oc-infinite-u}
\end{equation}
其中 $\mathbf P\succ\mathbf0$ 为常阵，满足\emph{连续代数 Riccati 方程}（CARE）
\begin{equation}
  \mathbf P\mathbf A+\mathbf A^\top\mathbf P-\mathbf P\mathbf B\mathbf R^{-1}\mathbf B^\top\mathbf P+\mathbf Q=\mathbf0.
  \label{eq:oc-care}
\end{equation}
\end{theorem}

\noindent 四点要紧（也是无限时域与有限时域的分水岭）：
\begin{enumerate}
  \item \textbf{要求能控}（有限时域不强求）：因控制区间扩到无穷，若不能控，则有发散方向使 $J=\infty$，无法比较优劣；能控性正是保证“代价可比”的条件。
  \item \textbf{终端项失效}：$t_f\to\infty$ 时 $\tfrac12\mathbf x^\top(t_f)\mathbf Q_f\mathbf x(t_f)$ 失去意义，即 $\mathbf Q_f=\mathbf0$。
  \item \textbf{定常闭环}：$\mathbf P$ 为常阵，闭环 $\dot{\mathbf x}=(\mathbf A-\mathbf B\mathbf K)\mathbf x$ 是线性定常系统。
  \item \textbf{闭环渐近稳定}：$\mathbf A-\mathbf B\mathbf K$ 特征值均具负实部，\emph{无论原系统 $\mathbf A$ 的特征值如何}。
\end{enumerate}

\begin{theorem}[闭环稳定：值函数即 Lyapunov 函数]\label{thm:oc-lqr-stable}
\cref{thm:oc-infinite-lqr} 的闭环渐近稳定。
\end{theorem}
\begin{proof}（展开导论 \cref{thm:ctrl-lqr-lyap} 的连续版，用 CARE）
取 $V(\mathbf x)=\mathbf x^\top\mathbf P\mathbf x\succ\mathbf0$（$\mathbf P\succ\mathbf0$）。沿闭环 $\dot{\mathbf x}=(\mathbf A-\mathbf B\mathbf R^{-1}\mathbf B^\top\mathbf P)\mathbf x$ 求导：
\[
  \dot V=\mathbf x^\top\big[(\mathbf A-\mathbf B\mathbf R^{-1}\mathbf B^\top\mathbf P)^\top\mathbf P+\mathbf P(\mathbf A-\mathbf B\mathbf R^{-1}\mathbf B^\top\mathbf P)\big]\mathbf x
  =\mathbf x^\top\big[(\mathbf A^\top\mathbf P+\mathbf P\mathbf A-\mathbf P\mathbf B\mathbf R^{-1}\mathbf B^\top\mathbf P)-\mathbf P\mathbf B\mathbf R^{-1}\mathbf B^\top\mathbf P\big]\mathbf x.
\]
由 CARE \cref{eq:oc-care}，括号内第一组 $=-\mathbf Q$，故 $\dot V=-\mathbf x^\top(\mathbf Q+\mathbf P\mathbf B\mathbf R^{-1}\mathbf B^\top\mathbf P)\mathbf x\le0$。$\mathbf Q\succ\mathbf0$ 时 $\dot V<0$ 负定，直接渐近稳定；$\mathbf Q\succeq\mathbf0$ 且 $\dot V$ 沿非平凡轨线不恒零时，由 LaSalle（\cref{ch:lyapunov-basics}）得渐近稳定。这把 LQR 闭环稳定性焊到了 \cref{ch:lyapunov-basics} 的 Lyapunov 方程 $\mathbf A^\top\mathbf P+\mathbf P\mathbf A=-\mathbf Q$ 上——LQR 是“自动构造 Lyapunov 函数”的机器。
\end{proof}

\begin{example}[二阶调节器：权重如何塑造极点]\label{ex:oc-lqr-second}
系统 $\dot{\mathbf x}=\begin{psmallmatrix}0&1\\0&0\end{psmallmatrix}\mathbf x+\begin{psmallmatrix}0\\1\end{psmallmatrix}u$（双积分），$J=\tfrac12\int_0^\infty(x_1^2+2bx_1x_2+ax_2^2+u^2)\,\mathrm dt$\footnote{控制项为 $+u^2$（刘豹勘误 M2-f：教材误印 $+u$，性能泛函须二次型，由解中 $\mathbf R=q_2=1$ 对应 $u^2$ 项可证）。}，即 $\mathbf Q=\begin{psmallmatrix}1&b\\b&a\end{psmallmatrix}$（需 $a-b^2>0$ 保正定）、$\mathbf R=1$。代数 Riccati 三式 $p_{12}^2=1,\ -p_{11}+p_{12}p_{22}-b=0,\ -2p_{12}+p_{22}^2-a=0$，取正定解 $p_{12}=1,\ p_{22}=\sqrt{a+2},\ p_{11}=\sqrt{a+2}-b$，故
\begin{equation}
  u^*=-x_1-\sqrt{a+2}\,x_2.
  \label{eq:oc-lqr-second-u}
\end{equation}
闭环极点 $s_{1,2}=-\tfrac{\sqrt{a+2}}{2}\pm\mathrm j\tfrac{\sqrt{2-a}}{2}$（$a<2$）。\textbf{权重塑造极点的直觉}：$a=0$（不罚 $x_2$）时极点为 $-\tfrac{\sqrt2}{2}\pm\mathrm j\tfrac{\sqrt2}{2}$，恰是经典控制里阻尼比 $\xi=0.707$ 的“最佳阻尼”二阶系统；$a$ 增大（越罚速度 $x_2$）极点趋向实轴、振荡减弱、响应变缓；$a>2$ 时过阻尼、无振荡。原系统不稳定（双重极点在原点），但 LQR 闭环渐近稳定——印证 \cref{thm:oc-lqr-stable}。
\end{example}

\begin{insight}
\cref{ex:oc-lqr-second} 把抽象的“选 $\mathbf Q,\mathbf R$”变成了可视的“挪极点”：调大 $\mathbf Q$ 对应分量（更罚误差）$\to$ 极点左移、响应更快但控制更猛；调大 $\mathbf R$（更省力）$\to$ 极点右移、响应更柔。LQR 因此可看作“自动极点配置”——你不必直接指定极点（如 \cref{ch:linear-synthesis} 的 Ackermann），而是通过物理意义明确的代价权重\emph{间接}塑造闭环。这正是 LQR 比手工极点配置更受工程青睐的根由，也是 \cref{ex:oc-static-lagrange} “种子”规律的最终果实。导论 \cref{sec:ctrl-lqr} 给的双推导回答“增益怎么算”，本例回答的是“权重怎么选、闭环长什么样”。
\end{insight}

\subsection{输出调节器与跟踪器}
\label{sec:oc-tracker}

\paragraph{输出调节器：罚输出而非全状态。}
若只要求\emph{输出} $\mathbf y=\mathbf C\mathbf x$ 接近平衡，代价改为 $J=\tfrac12\int[\mathbf y^\top\mathbf Q\mathbf y+\mathbf u^\top\mathbf R\mathbf u]\,\mathrm dt+\dots$。用 $\mathbf y=\mathbf C\mathbf x$ 代入，$\mathbf y^\top\mathbf Q\mathbf y=\mathbf x^\top(\mathbf C^\top\mathbf Q\mathbf C)\mathbf x$——\textbf{形式上就是把状态调节器的 $\mathbf Q$ 换成 $\mathbf C^\top\mathbf Q\mathbf C$}（系统能观时 $\mathbf C^\top\mathbf Q\mathbf C$ 仍半正定）。Riccati 与最优控制结构不变，只是 $\mathbf P(t)$ 不同。

\begin{pitfall}[“罚输出就该用输出反馈”]
直觉以为：既然代价罚的是输出 $\mathbf y$，最优控制就该是输出反馈 $\mathbf u=-\mathbf K_y\mathbf y$。\textbf{错}——最优控制仍是\emph{状态}反馈 $\mathbf u=-\mathbf K\mathbf x$。原因：状态矢量包含主宰过程未来演变的\emph{全部}信息，而输出只含部分信息；最优控制必须利用全部信息，故用 $\mathbf x$ 而非 $\mathbf y$ 反馈\cite{liubao2006modern}。若状态不可测，则需先用观测器重构 $\hat{\mathbf x}$（\cref{ch:linear-synthesis}），再 $\mathbf u=-\mathbf K\hat{\mathbf x}$——这就引出 LQG（LQR + Kalman 滤波）与分离原理，详见 \cref{ch:lqr-lqg-riccati} 与 \cref{ch:eskf}。
\end{pitfall}

\paragraph{跟踪器：让输出跟住参考。}
\emph{跟踪器}要使输出 $\mathbf y(t)$ 紧随期望输出 $\mathbf z(t)$，定义误差 $\mathbf e(t)=\mathbf z(t)-\mathbf C(t)\mathbf x(t)$，代价 $J=\tfrac12\int[\mathbf e^\top\mathbf Q\mathbf e+\mathbf u^\top\mathbf R\mathbf u]\,\mathrm dt+\tfrac12\mathbf e^\top(t_f)\mathbf Q_f\mathbf e(t_f)$。Hamilton 函数 $H=\tfrac12(\mathbf z-\mathbf C\mathbf x)^\top\mathbf Q(\mathbf z-\mathbf C\mathbf x)+\tfrac12\mathbf u^\top\mathbf R\mathbf u+\boldsymbol\lambda^\top(\mathbf A\mathbf x+\mathbf B\mathbf u)$。控制方程给 $\mathbf u^*=-\mathbf R^{-1}\mathbf B^\top\boldsymbol\lambda$，协态方程多出一个由 $\mathbf z$ 驱动的项。这逼使协态的扫描形式比调节器多一项：

\begin{theorem}[跟踪器最优控制]\label{thm:oc-tracker}
设扫描 ansatz $\boldsymbol\lambda(t)=\mathbf P(t)\mathbf x(t)-\mathbf g(t)$（比调节器的 $\boldsymbol\lambda=\mathbf P\mathbf x$ 多一个 $-\mathbf g$ 项），则最优控制
\begin{equation}
  \boxed{\mathbf u^*(t)=-\mathbf R^{-1}\mathbf B^\top\mathbf P(t)\mathbf x(t)+\mathbf R^{-1}\mathbf B^\top\mathbf g(t),}
  \label{eq:oc-tracker-u}
\end{equation}
其中 $\mathbf P(t)$ 满足与\emph{输出调节器相同}的 Riccati 方程 $\dot{\mathbf P}=-\mathbf P\mathbf A-\mathbf A^\top\mathbf P+\mathbf P\mathbf B\mathbf R^{-1}\mathbf B^\top\mathbf P-\mathbf C^\top\mathbf Q\mathbf C$，$\mathbf P(t_f)=\mathbf C^\top(t_f)\mathbf Q_f\mathbf C(t_f)$；前馈矢量 $\mathbf g(t)$ 满足\emph{伴随}方程
\begin{equation}
  \dot{\mathbf g}(t)=-[\mathbf A(t)-\mathbf B(t)\mathbf R^{-1}\mathbf B^\top(t)\mathbf P(t)]^\top\mathbf g(t)-\mathbf C^\top(t)\mathbf Q(t)\mathbf z(t),\quad\mathbf g(t_f)=\mathbf C^\top(t_f)\mathbf Q_f\mathbf z(t_f).
  \label{eq:oc-tracker-g}
\end{equation}
\end{theorem}

\begin{insight}
\textbf{反馈 $+$ 前馈：跟踪器的解剖学。}\; \cref{eq:oc-tracker-u} 第一项 $-\mathbf R^{-1}\mathbf B^\top\mathbf P\mathbf x$ 与调节器\emph{完全相同}——同一个 Riccati、同一个闭环矩阵 $\mathbf A-\mathbf B\mathbf R^{-1}\mathbf B^\top\mathbf P$、同样的极点。这意味着\emph{跟踪器的动态性能与参考 $\mathbf z(t)$ 无关}，由系统与权重唯一确定。差异全在第二项前馈 $\mathbf R^{-1}\mathbf B^\top\mathbf g$——$\mathbf g(t)$ 集中体现了“跟踪 vs 调节”的本质区别，把调节器变成了跟踪器。
\end{insight}

\begin{pitfall}[跟踪器需要参考的“全部未来值”]
\cref{eq:oc-tracker-g} 的 $\mathbf g(t)$ 方程要从\emph{终端}逆时间解（或先算 $\mathbf g(t_0)$ 再正向积分），其解
\[
  \mathbf g(t)=\boldsymbol\Psi^{-1}(t_f,t)\mathbf g(t_f)+\int_{t_0}^{t_f}\boldsymbol\Psi^{-1}(\tau,t)\mathbf C^\top(\tau)\mathbf Q(\tau)\mathbf z(\tau)\,\mathrm d\tau
\]
依赖参考 $\mathbf z(\tau)$ 在\emph{整个区间} $[t_0,t_f]$ 的值。因此\textbf{现时最优控制 $\mathbf u^*(t)$ 依赖参考的全部未来值}——这是因果性与最优性的张力：现时控制只能影响未来响应，而未来误差与 $\mathbf z$ 的未来值有关，故最优控制必须“预知”整条参考。实践中或用预估值代替，或把 $\mathbf z$ 当随机过程求期望最优。这也是机器人轨迹跟踪里“前瞻 / preview 控制”的理论根据：必须预先掌握期望轨迹 $\mathbf z(t)$ 的变化规律才能真正最优。
\end{pitfall}

\begin{example}[定常跟踪器：跟一个常值参考]\label{ex:oc-tracker-const}
系统 $\dot{\mathbf x}=\begin{psmallmatrix}0&1\\0&0\end{psmallmatrix}\mathbf x+\begin{psmallmatrix}0\\1\end{psmallmatrix}u,\ \mathbf y=(1,0)\mathbf x$，参考 $z(t)=r$（常值），$J=\tfrac12\int_0^\infty\{(x_1-z)^2+u^2\}\,\mathrm dt$。Riccati 解（同 \cref{ex:oc-lqr-second} 的 $a=0,b=0$）$\mathbf P=\begin{psmallmatrix}\sqrt2&1\\1&\sqrt2\end{psmallmatrix}$。前馈：取 $t_f\to\infty$、$\dot{\mathbf g}=\mathbf0$，由 \cref{eq:oc-tracker-g} 解得 $g_2=z=r,\ g_1=\sqrt2 r$。最优控制
\begin{equation}
  u^*=-x_1-\sqrt2\,x_2+g_2=-x_1-\sqrt2\,x_2+r.
  \label{eq:oc-tracker-const-u}
\end{equation}
即“调节器反馈 $-x_1-\sqrt2x_2$”加“常值前馈 $r$”——前馈把闭环平衡点从原点平移到 $x_1=r$，使输出稳态跟住参考。
\end{example}

\subsection{线性二次型次优控制：受限结构与输出反馈}
\label{sec:oc-lq-suboptimal}

\paragraph{动机：最优反馈要全状态，可状态未必都测得到。}
\cref{sec:oc-state-regulator,sec:oc-tracker} 的调节器、跟踪器，其最优律 $\mathbf u^*=-\mathbf K\mathbf x$ 一律是\emph{全状态}反馈——这天经地义，因为“最优”就该让反映系统内部运动的\emph{全部}信息参与组合\cite{liubao2006modern}。但工程现实是：并非每个状态分量都可测或易测（如 \cref{ex:oc-subopt} 里的速度 $x_2$ 无传感器）。出路有二：(a) 先用观测器重构 $\hat{\mathbf x}$ 再 $\mathbf u=-\mathbf K\hat{\mathbf x}$——这是 LQG + 分离原理的最优路线，留待 \cref{ch:lqr-lqg-riccati,ch:eskf}；(b) 干脆\emph{限定}控制器结构，只用现成的输出 $\mathbf y=\mathbf C\mathbf x$ 作线性反馈 $\mathbf u=-\mathbf K\mathbf y$。路线 (b) 所用信息不全，性能必不及最优，故称\emph{次优}（准最优）控制。\textbf{注意它与 \cref{sec:oc-tracker} 的输出调节器迥异}：那里性能指标虽提取输出信息，最优律 $\mathbf u^*=-\mathbf R^{-1}\mathbf B^\top\mathbf P\mathbf x=-\mathbf K\mathbf x$ 仍是全状态反馈；这里 $\mathbf u=-\mathbf K\mathbf y$ 是货真价实的\emph{输出}反馈。

\paragraph{受限结构 + Lyapunov：把 $J$ 写成 $\tfrac12\mathbf x^\top(0)\mathbf P\mathbf x(0)$。}
设完全能控、能观系统及二次型指标
\begin{equation}
  \dot{\mathbf x}=\mathbf A\mathbf x+\mathbf B\mathbf u,\quad\mathbf y=\mathbf C\mathbf x,\qquad
  J=\tfrac12\int_0^\infty[\mathbf x^\top\mathbf Q\mathbf x+\mathbf u^\top\mathbf R\mathbf u]\,\mathrm dt,
  \label{eq:oc-subopt-sys}
\end{equation}
$\mathbf Q\succeq\mathbf0$、$\mathbf R\succ\mathbf0$。规定控制由输出线性负反馈构成：
\begin{equation}
  \mathbf u(t)=-\mathbf K\mathbf y(t)=-\mathbf K\mathbf C\mathbf x(t),
  \label{eq:oc-subopt-u}
\end{equation}
则闭环 $\dot{\mathbf x}=(\mathbf A-\mathbf B\mathbf K\mathbf C)\mathbf x=\bar{\mathbf A}\mathbf x$。代入 $\mathbf u=-\mathbf K\mathbf C\mathbf x$，指标只剩状态二次型
\begin{equation}
  J=\tfrac12\int_0^\infty\mathbf x^\top\big[\mathbf Q+\mathbf C^\top\mathbf K^\top\mathbf R\mathbf K\mathbf C\big]\mathbf x\,\mathrm dt=\tfrac12\int_0^\infty\mathbf x^\top\bar{\mathbf Q}\mathbf x\,\mathrm dt,\quad\bar{\mathbf Q}=\mathbf Q+\mathbf C^\top\mathbf K^\top\mathbf R\mathbf K\mathbf C.
  \label{eq:oc-subopt-Qbar}
\end{equation}
设计任务遂归结为：\emph{选输出反馈阵 $\mathbf K$ 使 \cref{eq:oc-subopt-Qbar} 极小}。这恰好可用 \cref{ch:lyapunov-basics} 的“Lyapunov 函数即二次型指标”技巧。先假定 $\bar{\mathbf A}$ 的特征值全具负实部（闭环渐近稳定），取 $V(\mathbf x)=\mathbf x^\top\mathbf P\mathbf x$，则 $\dot V=\mathbf x^\top[\bar{\mathbf A}^\top\mathbf P+\mathbf P\bar{\mathbf A}]\mathbf x$。令 $\mathbf P$ 满足 Lyapunov 方程（即 \cref{eq:lyap-eq} 用于闭环阵 $\bar{\mathbf A}$）
\begin{equation}
  \bar{\mathbf A}^\top\mathbf P+\mathbf P\bar{\mathbf A}=-\bar{\mathbf Q},
  \label{eq:oc-subopt-lyap}
\end{equation}
便得 $\dot V=-\mathbf x^\top\bar{\mathbf Q}\mathbf x$，即 $\mathbf x^\top\bar{\mathbf Q}\mathbf x=-\frac{\mathrm d}{\mathrm dt}(\mathbf x^\top\mathbf P\mathbf x)$。代回 \cref{eq:oc-subopt-Qbar} 并用 $\bar{\mathbf A}$ 稳定 $\Rightarrow\mathbf x(\infty)\to\mathbf0$：
\begin{equation}
  J=\tfrac12\int_0^\infty\mathbf x^\top\bar{\mathbf Q}\mathbf x\,\mathrm dt=-\tfrac12\mathbf x^\top\mathbf P\mathbf x\Big|_0^\infty=\boxed{\tfrac12\mathbf x^\top(0)\mathbf P\mathbf x(0).}
  \label{eq:oc-subopt-J}
\end{equation}

\begin{insight}
\textbf{同一副面孔，两种来历。}\; \cref{eq:oc-subopt-J} 与最优调节器的 $J^*=\tfrac12\mathbf x^\top(0)\mathbf P\mathbf x(0)$（\cref{eq:oc-lqr-Jstar}）\emph{形式完全相同}，但内涵不同：那里 $\mathbf P$ 解\emph{Riccati} 方程 \cref{eq:oc-care}、对应全状态最优；这里 $\mathbf P$ 解\emph{Lyapunov} 方程 \cref{eq:oc-subopt-lyap}、对应受限结构次优。由于输出反馈 $\mathbf u=-\mathbf K\mathbf y$ 只是全体镇定控制律的一个\emph{真子类}，故对任意 $\mathbf K$ 都有\textbf{次优界} $J_{\text{次优}}(\mathbf K)\ge J^*_{\text{LQR}}$；等号仅当最优全状态律恰能由输出反馈实现时成立。这把“次优代价比最优贵多少”量化成了 Lyapunov 解与 Riccati 解之差 $\tfrac12\mathbf x^\top(0)(\mathbf P_{\text{Lyap}}-\mathbf P_{\text{Riccati}})\mathbf x(0)\ge0$。
\end{insight}

\paragraph{求 $\mathbf K$：梯度法与对初值的依赖。}
反馈阵 $\mathbf K$ 不能从 Lyapunov 方程 $(\mathbf A-\mathbf B\mathbf K\mathbf C)^\top\mathbf P+\mathbf P(\mathbf A-\mathbf B\mathbf K\mathbf C)=-\mathbf Q-\mathbf C^\top\mathbf K^\top\mathbf R\mathbf K\mathbf C$ 直接解出，因其中 $\mathbf P$ 与 $\mathbf K$ \emph{皆未知}。一个可行办法是\emph{梯度速降}：由该方程解出以 $\mathbf K$ 表示的 $\mathbf P[\mathbf K]$，代入 \cref{eq:oc-subopt-J} 得 $J(\mathbf K)$，再令
\begin{equation}
  \frac{\partial J(\mathbf K)}{\partial\mathbf K}=\mathbf0
  \label{eq:oc-subopt-grad}
\end{equation}
解出使 $J\to\min$ 的 $\mathbf K$。\textbf{一个本质性的代价}：如此求得的最优 $\mathbf K$ \emph{与初始条件 $\mathbf x(0)$ 有关}（因 \cref{eq:oc-subopt-J} 显含 $\mathbf x(0)$）——这与全状态 LQR 增益 $\mathbf K=\mathbf R^{-1}\mathbf B^\top\mathbf P$ \emph{不依赖} $\mathbf x(0)$ 形成鲜明对照，是受限结构次优控制的固有短板。工程上或对一族典型 $\mathbf x(0)$ 取平均（视 $\mathbf x(0)$ 为随机、最小化 $\mathbb E[J]$），或接受“为某一标称初值调优”。

\begin{example}[输出反馈次优控制：$x_2$ 不可测]\label{ex:oc-subopt}
受控系统 $\dot{\mathbf x}=\begin{psmallmatrix}0&1\\0&-1\end{psmallmatrix}\mathbf x+\begin{psmallmatrix}0\\1\end{psmallmatrix}u$，$\mathbf y=(1,0)\mathbf x$（只测得位置 $x_1$、速度 $x_2$ 不可测），$\mathbf Q=\mathbf I_2$、$\mathbf R=1$。采用输出反馈 $u=-k\,y=-k\,x_1$ 求标量增益 $k$ 使 $J\to\min$。闭环阵 $\bar{\mathbf A}=\mathbf A-\mathbf B k\mathbf C=\begin{psmallmatrix}0&1\\-k&-1\end{psmallmatrix}$；权阵 $\bar{\mathbf Q}=\mathbf Q+\mathbf C^\top k\mathbf R k\mathbf C=\begin{psmallmatrix}1+k^2&0\\0&1\end{psmallmatrix}$。解 Lyapunov 方程 $\bar{\mathbf A}^\top\mathbf P+\mathbf P\bar{\mathbf A}=-\bar{\mathbf Q}$（$\mathbf P=\begin{psmallmatrix}p_{11}&p_{12}\\p_{12}&p_{22}\end{psmallmatrix}$）得三式 $2kp_{12}=1+k^2$、$p_{11}-p_{12}-kp_{22}=0$、$2p_{12}-2p_{22}=-1$，联立
\[
  p_{12}=\frac{1+k^2}{2k},\quad p_{22}=\frac{1+k+k^2}{2k},\quad p_{11}=\frac{1+k+2k^2+k^3}{2k}.
\]
取初值 $\mathbf x(0)=(1,0)^\top$，则 $J(k)=\tfrac12\mathbf x^\top(0)\mathbf P\mathbf x(0)=\tfrac12p_{11}=\dfrac{1+k+2k^2+k^3}{4k}$。令 $\dfrac{\partial J}{\partial k}=\dfrac14\big(2k+2-k^{-2}\big)=0$，即 $2k^3+2k^2-1=0$，解得
\begin{equation}
  k^*\approx0.57,\qquad J(k^*)\approx1.05.
  \label{eq:oc-subopt-result}
\end{equation}
此即 $x_2$ 不可测时、用单一位置反馈所能达到的最好（次优）性能；其代价 $1.05$ 必不小于同 $\mathbf Q,\mathbf R$ 下全状态 LQR 的最优值——差额就是“放弃速度信息”的定量惩罚。若换一个 $\mathbf x(0)$，最优 $k^*$ 亦随之改变，再次印证 $\mathbf K$ 对初值的依赖。
\end{example}

\begin{pitfall}[次优控制别当成“另一种最优”]
输出反馈次优控制读起来与 LQR 很像（同样的 $\tfrac12\mathbf x^\top(0)\mathbf P\mathbf x(0)$、同样解一个矩阵方程），易误以为它是“输出版的最优控制”。\textbf{两点须分清}：其一，它的 $\mathbf P$ 解 Lyapunov 方程而非 Riccati，给出的是\emph{受限结构内}的最好、\emph{不是}全局最优（\cref{thm:oc-infinite-lqr} 才是）；其二，它的增益依赖 $\mathbf x(0)$，不是一劳永逸的反馈律。真正“最优地”利用不完全量测信息的，是先估计再控制的 LQG（\cref{ch:lqr-lqg-riccati}）——由分离原理，LQG 用 Kalman 滤波 \cref{ch:eskf} 重构状态、再套全状态 LQR 增益，其增益\emph{不}依赖 $\mathbf x(0)$。次优输出反馈胜在结构简单、无需观测器，适合低阶或对成本敏感的场合。
\end{pitfall}

\subsection{与导论的离散 LQR 一个口径}
\label{sec:oc-lqr-discrete-preview}

\paragraph{三条路、一个 Riccati。}
至此，LQR 的 Riccati 方程已被本书三条独立路径导出：导论用\emph{动态规划}（\cref{thm:ctrl-lqr-dp}，离散）与\emph{HJB}（\cref{eq:ctrl-care}，连续）两条，本章用\emph{极小值原理}（协态扫描，\cref{thm:oc-rde}）第三条。三者必然一致——因为它们都是同一最优控制问题的必要（对线性二次型也充分）条件，只是从 Hamilton 框架的不同侧面切入。离散侧，\cref{ex:oc-dp-scalar} 已用标量 DP 例题点醒“最优反馈增益是状态系数”；一般离散 LQR 的完整归纳推导与离散 Riccati 递推 \cref{eq:ctrl-dlqr-P}、DARE \cref{eq:ctrl-dare} 见导论 \cref{thm:ctrl-lqr-dp,thm:ctrl-lqr-converge}，本章不重推。

\begin{note}
\textbf{“最优 $\ne$ 稳定”一句提醒.}\; 导论 \cref{ex:ctrl-kalman}（Kalman 反例）警示：\emph{有限}时域最优未必给稳定闭环，唯\emph{无限}时域 LQR 才有稳定性保证（\cref{thm:oc-lqr-stable}）——因为它把无穷远未来的代价也算进去、不为短期收益牺牲长期稳定。本章的有限时域调节器（\cref{sec:oc-state-regulator}）正是“有限时域”，其首段增益 $\mathbf K(t_0)$ 单独作定常律时可能不稳定；用整条时变 $\mathbf K(t)$ 才最优。这与 \cref{ex:oc-lqr-scalar} 的“$t_f$ 越大、$P(t)$ 越接近稳态”图景互为表里。MPC 用终端代价 + 终端约束“伪造”无限时域来修复稳定性，详见 \cref{sec:ctrl-mpc}。
\end{note}

% ════════════════════════════════════════════════════════════════
\section{通向高级：本章在最优控制版图中的位置}
\label{sec:oc-outlook}

本章把最优控制“第一次系统讲清”：四件工具——变分法、极小值原理、动态规划、LQR——在 Hamilton 框架下统一，每件都建立了物理直觉与几何图解。这层地基之上，A--Q 高级章把每件工具严格化或推广到前沿：

\begin{itemize}
  \item \textbf{变分 / 极小值原理} $\to$ \cref{ch:opt-variational-pmp}：本章的极小值原理 \cref{thm:oc-pmp} 在那里被严格化为\emph{针状变分}（needle variation）证明、并处理\emph{奇异弧}（$H$ 对 $\mathbf u$ 退化、Bang--Bang 失效的边缘情形）；本章 \cref{ex:oc-pmp-linear} 的切换结构在那里获得完整的奇异/正则分类。
  \item \textbf{动态规划 / HJB} $\to$ \cref{ch:dp-hjb}：本章 \cref{thm:oc-hjb} 的 HJB 在那里以\emph{黏性解}（viscosity solution）严格处理 $J^*$ 不可微的情形（如 Bang--Bang 开关曲线处），并给出\emph{值迭代 / 策略迭代}（VI/PI）的收敛理论与数值解法。
  \item \textbf{LQR / Riccati} $\to$ \cref{ch:lqr-lqg-riccati}：本章 \cref{thm:oc-infinite-lqr} 的 CARE \cref{eq:oc-care} 在那里获得\emph{存在唯一性}证明（可镇定 + 可检测 $\Rightarrow$ 唯一正定镇定解）、\emph{分离原理}（LQG = LQR + Kalman，确定性等价）完整证明、\emph{回差不等式}与到 $H_\infty$ / 博弈型 Riccati 的推广；本章的输出调节器/跟踪器在那里接上 LQG 与 \cref{ch:eskf} 的“估计--控制对偶”。
  \item \textbf{Lyapunov 接口} $\to$ \cref{ch:lyapunov-basics} 与高级稳定性章：本章 \cref{thm:oc-lqr-stable} 用“值函数即 Lyapunov 函数”证 LQR 闭环稳定，把最优控制与 Lyapunov 理论焊接；HJB 的最优值函数是更一般非线性系统“控制 Lyapunov 函数”的来源（接 \cref{ch:control} 的 CLF 思想）。
\end{itemize}

\noindent 在机器人学语境，本章工具是后续应用章的理论上游：规划侧的 iLQR/DDP（\cref{sec:ctrl-ilqr-ddp}，HJB/DP 的局部二次近似）、MPC（\cref{sec:ctrl-mpc}，有限时域最优的滚动实现）、足式 WBC/凸 MPC（\cref{ch:control} 足式部分，LQR/QP 在浮动基座动力学 $\mathbf M(\mathbf q)\ddot{\mathbf q}+\mathbf C\dot{\mathbf q}+\mathbf g=\boldsymbol\tau$ 上的实例化）皆以此为根。读完本章，你应能在“变分 $\to$ 极小值原理 $\to$ HJB $\to$ LQR”这条链上自由游走，并知道每个结论该到哪一章去求严格证明或前沿推广。

% ════════════════════════════════════════════════════════════════
\section*{本章小结}

\begin{table}[H]\centering\small
\caption{符号表}
\begin{tabular}{ll}
\toprule
符号 & 含义 \\
\midrule
$\mathbf{x}(t),\mathbf{u}(t)$ & 状态矢量（$\mathbb R^n$）/ 控制矢量（$\mathbb R^r$） \\
$J,\ L,\ \Phi$ & 性能泛函 / 代价密度（Lagrange 型）/ 终端代价（Mayer 型） \\
$\boldsymbol\lambda(t)$ & 协态 / 伴随矢量（$=\partial J^*/\partial\mathbf x$） \\
$H=L+\boldsymbol\lambda^\top\mathbf f$ & Hamilton 函数 \\
$\delta J,\ \delta y$ & 泛函的变分 / 宗量的变分 \\
$\mathbf Q\succeq\mathbf0,\ \mathbf R\succ\mathbf0,\ \mathbf Q_f\succeq\mathbf0$ & 状态权 / 控制权 / 终端权（二次型代价） \\
$\mathbf P(t),\ \mathbf P$ & Riccati 矩阵（cost-to-go 权重，时变 / 常阵） \\
$\mathbf K=\mathbf R^{-1}\mathbf B^\top\mathbf P$ & LQR 最优反馈增益（$\mathbf u^*=-\mathbf K\mathbf x$） \\
$\mathbf g(t)$ & 跟踪器前馈矢量 \\
$\gamma,\ h(x_1,x_2)$ & Bang--Bang 开关曲线及其切换函数 \\
$J^*[\mathbf x,t]$ & 最优值函数 / cost-to-go \\
$\operatorname{SGN},\operatorname{sgn}$ & 矢量 / 标量符号函数 \\
\bottomrule
\end{tabular}
\end{table}

\begin{table}[H]\centering\small
\caption{定理与方程速查表}
\begin{tabular}{lll}
\toprule
名称 & 标签 & 一句话 \\
\midrule
泛函极值必要条件 & \cref{thm:oc-variation-zero} & $\delta J=0$（“梯度为零”的无穷维版） \\
Euler--Lagrange 方程 & \cref{thm:oc-euler} & $L_x-\tfrac{\mathrm d}{\mathrm dt}L_{\dot x}=0$ \\
横截条件 & \cref{eq:oc-transversality-free,eq:oc-transversality-target} & 自由 / 可变端点补足边界 \\
Hamilton 正则方程 & \cref{thm:oc-canonical} & $\dot{\mathbf x}=\partial_{\boldsymbol\lambda}H,\ \dot{\boldsymbol\lambda}=-\partial_{\mathbf x}H,\ \partial_{\mathbf u}H=0$ \\
Pontryagin 极小值原理 & \cref{thm:oc-pmp} & 控制受界时 $\partial_{\mathbf u}H=0\to\min_{\mathbf u\in U}H$ \\
Bang--Bang 控制 & \cref{thm:oc-bangbang} & $\mathbf u^*=-\operatorname{SGN}(\mathbf B^\top\boldsymbol\lambda)$，开关曲线 \\
Bellman 最优性原理 & \cref{thm:oc-bellman-principle} & 最优轨线的后段也最优 \\
HJB 方程 & \cref{thm:oc-hjb} & $-\partial_t J^*=\min_{\mathbf u}H$，$\boldsymbol\lambda=\partial_{\mathbf x}J^*$ \\
Riccati 微分方程（RDE）& \cref{thm:oc-rde} & 有限时域调节器，$\mathbf P(t_f)=\mathbf Q_f$ \\
无限时域 LQR / CARE & \cref{thm:oc-infinite-lqr} & 能控 $\Rightarrow$ 常增益 $\mathbf u^*=-\mathbf K\mathbf x$，闭环稳定 \\
跟踪器 & \cref{thm:oc-tracker} & 调节器反馈 $+$ 前馈 $\mathbf g(t)$（依赖参考全程） \\
离散正则方程 & \cref{thm:oc-disc-canonical} & $\mathbf x(k{+}1)=\partial_{\boldsymbol\lambda}H_k,\ \boldsymbol\lambda(k)=\partial_{\mathbf x}H_k,\ \partial_{\mathbf u}H_k=0$ \\
LQ 次优（输出反馈）& \cref{sec:oc-lq-suboptimal} & $\mathbf u=-\mathbf K\mathbf y$，$\mathbf P$ 解 Lyapunov 方程，$J\ge J^*_{\text{LQR}}$ \\
\bottomrule
\end{tabular}
\end{table}

\begin{table}[H]\centering\small
\caption{四件工具的适用边界（知识点总表）}
\begin{tabular}{llll}
\toprule
工具 & 适用情形 & 给出 & 局限 \\
\midrule
变分法 & 控制\emph{无界}、$L$ 光滑 & 局部必要条件（开环） & 控制受界时失效 \\
极小值原理 & 控制\emph{受界} $\mathbf u\in U$ & $\min_{\mathbf u\in U}H$（开环） & 仅必要（线性系统也充分） \\
动态规划 & 受界 / 不受界 & 全局\emph{反馈}解、HJB & $J^*$ 须可微，维数灾难 \\
LQR & 线性 $+$ 二次代价 & 闭式线性反馈 $-\mathbf K\mathbf x$ & 仅线性二次型 \\
\bottomrule
\end{tabular}
\end{table}

\section*{常见误解汇总}
\begin{enumerate}
  \item \textbf{“最优一定稳定。”}\; 错。有限时域最优未必稳定（导论 Kalman 反例 \cref{ex:ctrl-kalman}），唯无限时域 LQR 有稳定性保证（\cref{thm:oc-lqr-stable}）。
  \item \textbf{“控制受界还能用 $\partial H/\partial\mathbf u=0$。”}\; 错。受界时平稳条件失效，须用极小值原理的全局极小 \cref{eq:oc-pmp-min}（\cref{ex:oc-pmp-linear}）。
  \item \textbf{“罚输出就用输出反馈。”}\; 错。最优控制仍是状态反馈 $-\mathbf K\mathbf x$，因状态含全部未来信息（\cref{sec:oc-tracker} 陷阱）。
  \item \textbf{“跟踪器现时控制只需现时参考。”}\; 错。$\mathbf u^*(t)$ 依赖参考 $\mathbf z$ 的\emph{全部未来值}（\cref{eq:oc-tracker-g} 陷阱）。
  \item \textbf{“变分法、极小值原理、动态规划是三套无关方法。”}\; 错。三者经 $\boldsymbol\lambda=\partial J^*/\partial\mathbf x$ 统一于 Hamilton 框架，互为表里（\cref{thm:oc-hjb} 邻近 insight）。
  \item \textbf{“输出反馈次优控制就是输出版的最优控制。”}\; 错。其 $\mathbf P$ 解 Lyapunov 方程（非 Riccati），只是受限结构内最好、且增益依赖 $\mathbf x(0)$；真正最优地利用不完全量测的是 LQG（\cref{sec:oc-lq-suboptimal} 陷阱）。
  \item \textbf{“连续极小值原理可原样搬到离散时间。”}\; 错。离散“$\arg\min_{\mathbf u\in U}H_k$”需附加凸性条件才必要，无界时只保平稳形式 $\partial H_k/\partial\mathbf u=\mathbf0$（\cref{thm:oc-disc-canonical} 后陷阱）。
\end{enumerate}

\section*{延伸阅读}
\begin{itemize}
  \item \textbf{主源（物理直觉 + 例题密集）}：刘豹《现代控制理论》第3版\cite{liubao2006modern} 第6章——LP 引入、二次型物理意义、Bang--Bang 双积分相平面、DP/HJB、LQR 调节器/跟踪器的完整中文叙事，本章正文主线的出处。
  \item \textbf{第二信源（互校 + 四端点变分）}：张嗣瀛 / 高立群《现代控制理论》第2版\cite{zhang2017modern} 第7章——四类端点条件的变分解法、横截条件、LQ 从 Hamilton 到 Riccati 的另一条解析路线，与刘豹互为人审 ground-truth。
  \item \textbf{经典溯源}：Pontryagin 等《最优过程的数学理论》\cite{pontryagin1962mathematical}（极小/极大值原理原典）；Bellman《动态规划》\cite{bellman1957dynamic}（DP 与最优性原理原典）；Bryson \& Ho《Applied Optimal Control》\cite{bryson1975applied}（变分 / 协态 / 估计统一处理）；Kirk《Optimal Control Theory》\cite{kirk2004optimal}（教学经典）；Anderson \& Moore《Optimal Control: Linear Quadratic Methods》\cite{anderson2007optimal}（LQR/LQG 权威）；Lewis 等《Optimal Control》\cite{lewis2012optimal}。
  \item \textbf{机器人语境}：Tedrake《Underactuated Robotics》\cite{tedrake_underactuated}（LQR/iLQR/HJB 在欠驱动机器人上的现代讲法，与导论 \cref{sec:ctrl-lqr,sec:ctrl-ilqr-ddp} 共享）。
\end{itemize}

\section*{本章与后续章节的关系}
\begin{table}[H]\centering\small
\begin{tabular}{lll}
\toprule
后续章节 & 关系 & 本章铺垫 \\
\midrule
\cref{ch:opt-variational-pmp}（PMP）& 严格化 & 极小值原理 \cref{thm:oc-pmp}、奇异弧（\cref{ex:oc-pmp-linear} 切换结构） \\
\cref{ch:dp-hjb}（DP/HJB）& 严格化 & HJB \cref{thm:oc-hjb}、黏性解 / VI-PI（$J^*$ 不可微情形） \\
\cref{ch:lqr-lqg-riccati}（LQR/Riccati）& 严格化 & CARE \cref{eq:oc-care} 存在唯一、分离原理、$H_\infty$ \\
\cref{ch:lyapunov-basics}（Lyapunov）& 接口 & “值函数即 Lyapunov”\cref{thm:oc-lqr-stable} \\
\cref{sec:ctrl-ilqr-ddp}（iLQR/DDP）& 下游应用 & DP/HJB 的局部二次近似 \\
\cref{sec:ctrl-mpc}（MPC）& 下游应用 & 有限时域最优的滚动实现、“伪造”无限时域 \\
\cref{ch:eskf}（状态估计）& 横向对偶 & 估计--控制对偶（LQG 的另一半） \\
\bottomrule
\end{tabular}
\end{table}

\section*{故障排查手册}
\begin{table}[H]\centering\small
\begin{tabular}{p{0.36\linewidth}p{0.58\linewidth}}
\toprule
症状 & 可能原因与对策 \\
\midrule
Euler--Lagrange 解欠定 / 多解 & 漏了横截条件：每个自由端点须配一个横截条件 \cref{eq:oc-transversality-free,eq:oc-transversality-target} 补足边界。 \\
变分法求不出受界最优控制 & 控制受界时 $\partial H/\partial\mathbf u=0$ 失效，改用极小值原理 $\min_{\mathbf u\in U}H$（\cref{thm:oc-pmp}）；若 $H$ 对 $\mathbf u$ 线性，最优控制取边界值（Bang--Bang）。 \\
Bang--Bang 切换次数算不对 & 切换次数上界 $n-1$ 仅当 $\mathbf A$ 特征值全为实数；有复特征值时无上界（\cref{thm:oc-bangbang} 后注）。 \\
LQR 闭环不稳定 & 检查是否\emph{有限}时域单段增益当常律用了（导论 Kalman 反例 \cref{ex:ctrl-kalman}）；无限时域须系统\emph{能控}、$\mathbf R\succ\mathbf0$、$(\mathbf A,\mathbf Q^{1/2})$ 可检测。 \\
跟踪器稳态有偏差 & 检查前馈 $\mathbf g(t)$ 是否正确（\cref{eq:oc-tracker-g}）；终端 $t_f$ 附近误差变大常因 $\mathbf Q_f=\mathbf0\Rightarrow\mathbf g(t_f)=\mathbf0\Rightarrow\mathbf u(t_f)=0$。 \\
HJB 解不出 & $H$ 对 $\mathbf u$ 非二次 / $J^*$ 不光滑时一般无解析解，需数值（值迭代）；严格的黏性解理论见 \cref{ch:dp-hjb}。 \\
权重 $\mathbf Q,\mathbf R$ 怎么选 & 用物理意义（\cref{sec:oc-quadratic-meaning}）：调大 $\mathbf Q$ 分量更罚该误差、极点左移更快；调大 $\mathbf R$ 更省力、极点右移更柔（\cref{ex:oc-lqr-second}）。 \\
状态不全可测、仍要状态反馈 & 两条路：(a) 观测器重构 $\hat{\mathbf x}$ 走 LQG（最优、增益不依赖 $\mathbf x(0)$，\cref{ch:lqr-lqg-riccati}）；(b) 限定输出反馈 $\mathbf u=-\mathbf K\mathbf y$ 走 LQ 次优（\cref{sec:oc-lq-suboptimal}，$\mathbf K$ 依赖 $\mathbf x(0)$）。 \\
离散问题受界控制求不出 & 离散 $\arg\min_{\mathbf u\in U}H_k$ 需凸性才必要；无界时用平稳形式 $\partial H_k/\partial\mathbf u=\mathbf0$ 与离散正则方程 \cref{thm:oc-disc-canonical}，严格论见 \cref{ch:opt-variational-pmp}。 \\
\bottomrule
\end{tabular}
\end{table}
